% !TEX root = appendix.tex
% Single-file appendix source.
% Compile this file directly for the standalone supplement, or define
% \ICLRIncludeAppendix before loading main.tex for the integrated submission.
\ifdefined\ICLRIncludeAppendix
\def\ICLRFinishAppendix{\endinput}
\else
\pdfminorversion=7
\documentclass{article}
\usepackage{iclr2027_conference,times}
%%%%% NEW MATH DEFINITIONS %%%%%

\usepackage{amsmath,amsfonts,bm}

% Mark sections of captions for referring to divisions of figures
\newcommand{\figleft}{{\em (Left)}}
\newcommand{\figcenter}{{\em (Center)}}
\newcommand{\figright}{{\em (Right)}}
\newcommand{\figtop}{{\em (Top)}}
\newcommand{\figbottom}{{\em (Bottom)}}
\newcommand{\captiona}{{\em (a)}}
\newcommand{\captionb}{{\em (b)}}
\newcommand{\captionc}{{\em (c)}}
\newcommand{\captiond}{{\em (d)}}

% Highlight a newly defined term
\newcommand{\newterm}[1]{{\bf #1}}


% Figure reference, lower-case.
\def\figref#1{figure~\ref{#1}}
% Figure reference, capital. For start of sentence
\def\Figref#1{Figure~\ref{#1}}
\def\twofigref#1#2{figures \ref{#1} and \ref{#2}}
\def\quadfigref#1#2#3#4{figures \ref{#1}, \ref{#2}, \ref{#3} and \ref{#4}}
% Section reference, lower-case.
\def\secref#1{section~\ref{#1}}
% Section reference, capital.
\def\Secref#1{Section~\ref{#1}}
% Reference to two sections.
\def\twosecrefs#1#2{sections \ref{#1} and \ref{#2}}
% Reference to three sections.
\def\secrefs#1#2#3{sections \ref{#1}, \ref{#2} and \ref{#3}}
% Reference to an equation, lower-case.
\def\eqref#1{equation~\ref{#1}}
% Reference to an equation, upper case
\def\Eqref#1{Equation~\ref{#1}}
% A raw reference to an equation---avoid using if possible
\def\plaineqref#1{\ref{#1}}
% Reference to a chapter, lower-case.
\def\chapref#1{chapter~\ref{#1}}
% Reference to an equation, upper case.
\def\Chapref#1{Chapter~\ref{#1}}
% Reference to a range of chapters
\def\rangechapref#1#2{chapters\ref{#1}--\ref{#2}}
% Reference to an algorithm, lower-case.
\def\algref#1{algorithm~\ref{#1}}
% Reference to an algorithm, upper case.
\def\Algref#1{Algorithm~\ref{#1}}
\def\twoalgref#1#2{algorithms \ref{#1} and \ref{#2}}
\def\Twoalgref#1#2{Algorithms \ref{#1} and \ref{#2}}
% Reference to a part, lower case
\def\partref#1{part~\ref{#1}}
% Reference to a part, upper case
\def\Partref#1{Part~\ref{#1}}
\def\twopartref#1#2{parts \ref{#1} and \ref{#2}}

\def\ceil#1{\lceil #1 \rceil}
\def\floor#1{\lfloor #1 \rfloor}
\def\1{\bm{1}}
\newcommand{\train}{\mathcal{D}}
\newcommand{\valid}{\mathcal{D_{\mathrm{valid}}}}
\newcommand{\test}{\mathcal{D_{\mathrm{test}}}}

\def\eps{{\epsilon}}


% Random variables
\def\reta{{\textnormal{$\eta$}}}
\def\ra{{\textnormal{a}}}
\def\rb{{\textnormal{b}}}
\def\rc{{\textnormal{c}}}
\def\rd{{\textnormal{d}}}
\def\re{{\textnormal{e}}}
\def\rf{{\textnormal{f}}}
\def\rg{{\textnormal{g}}}
\def\rh{{\textnormal{h}}}
\def\ri{{\textnormal{i}}}
\def\rj{{\textnormal{j}}}
\def\rk{{\textnormal{k}}}
\def\rl{{\textnormal{l}}}
% rm is already a command, just don't name any random variables m
\def\rn{{\textnormal{n}}}
\def\ro{{\textnormal{o}}}
\def\rp{{\textnormal{p}}}
\def\rq{{\textnormal{q}}}
\def\rr{{\textnormal{r}}}
\def\rs{{\textnormal{s}}}
\def\rt{{\textnormal{t}}}
\def\ru{{\textnormal{u}}}
\def\rv{{\textnormal{v}}}
\def\rw{{\textnormal{w}}}
\def\rx{{\textnormal{x}}}
\def\ry{{\textnormal{y}}}
\def\rz{{\textnormal{z}}}

% Random vectors
\def\rvepsilon{{\mathbf{\epsilon}}}
\def\rvtheta{{\mathbf{\theta}}}
\def\rva{{\mathbf{a}}}
\def\rvb{{\mathbf{b}}}
\def\rvc{{\mathbf{c}}}
\def\rvd{{\mathbf{d}}}
\def\rve{{\mathbf{e}}}
\def\rvf{{\mathbf{f}}}
\def\rvg{{\mathbf{g}}}
\def\rvh{{\mathbf{h}}}
\def\rvu{{\mathbf{i}}}
\def\rvj{{\mathbf{j}}}
\def\rvk{{\mathbf{k}}}
\def\rvl{{\mathbf{l}}}
\def\rvm{{\mathbf{m}}}
\def\rvn{{\mathbf{n}}}
\def\rvo{{\mathbf{o}}}
\def\rvp{{\mathbf{p}}}
\def\rvq{{\mathbf{q}}}
\def\rvr{{\mathbf{r}}}
\def\rvs{{\mathbf{s}}}
\def\rvt{{\mathbf{t}}}
\def\rvu{{\mathbf{u}}}
\def\rvv{{\mathbf{v}}}
\def\rvw{{\mathbf{w}}}
\def\rvx{{\mathbf{x}}}
\def\rvy{{\mathbf{y}}}
\def\rvz{{\mathbf{z}}}

% Elements of random vectors
\def\erva{{\textnormal{a}}}
\def\ervb{{\textnormal{b}}}
\def\ervc{{\textnormal{c}}}
\def\ervd{{\textnormal{d}}}
\def\erve{{\textnormal{e}}}
\def\ervf{{\textnormal{f}}}
\def\ervg{{\textnormal{g}}}
\def\ervh{{\textnormal{h}}}
\def\ervi{{\textnormal{i}}}
\def\ervj{{\textnormal{j}}}
\def\ervk{{\textnormal{k}}}
\def\ervl{{\textnormal{l}}}
\def\ervm{{\textnormal{m}}}
\def\ervn{{\textnormal{n}}}
\def\ervo{{\textnormal{o}}}
\def\ervp{{\textnormal{p}}}
\def\ervq{{\textnormal{q}}}
\def\ervr{{\textnormal{r}}}
\def\ervs{{\textnormal{s}}}
\def\ervt{{\textnormal{t}}}
\def\ervu{{\textnormal{u}}}
\def\ervv{{\textnormal{v}}}
\def\ervw{{\textnormal{w}}}
\def\ervx{{\textnormal{x}}}
\def\ervy{{\textnormal{y}}}
\def\ervz{{\textnormal{z}}}

% Random matrices
\def\rmA{{\mathbf{A}}}
\def\rmB{{\mathbf{B}}}
\def\rmC{{\mathbf{C}}}
\def\rmD{{\mathbf{D}}}
\def\rmE{{\mathbf{E}}}
\def\rmF{{\mathbf{F}}}
\def\rmG{{\mathbf{G}}}
\def\rmH{{\mathbf{H}}}
\def\rmI{{\mathbf{I}}}
\def\rmJ{{\mathbf{J}}}
\def\rmK{{\mathbf{K}}}
\def\rmL{{\mathbf{L}}}
\def\rmM{{\mathbf{M}}}
\def\rmN{{\mathbf{N}}}
\def\rmO{{\mathbf{O}}}
\def\rmP{{\mathbf{P}}}
\def\rmQ{{\mathbf{Q}}}
\def\rmR{{\mathbf{R}}}
\def\rmS{{\mathbf{S}}}
\def\rmT{{\mathbf{T}}}
\def\rmU{{\mathbf{U}}}
\def\rmV{{\mathbf{V}}}
\def\rmW{{\mathbf{W}}}
\def\rmX{{\mathbf{X}}}
\def\rmY{{\mathbf{Y}}}
\def\rmZ{{\mathbf{Z}}}

% Elements of random matrices
\def\ermA{{\textnormal{A}}}
\def\ermB{{\textnormal{B}}}
\def\ermC{{\textnormal{C}}}
\def\ermD{{\textnormal{D}}}
\def\ermE{{\textnormal{E}}}
\def\ermF{{\textnormal{F}}}
\def\ermG{{\textnormal{G}}}
\def\ermH{{\textnormal{H}}}
\def\ermI{{\textnormal{I}}}
\def\ermJ{{\textnormal{J}}}
\def\ermK{{\textnormal{K}}}
\def\ermL{{\textnormal{L}}}
\def\ermM{{\textnormal{M}}}
\def\ermN{{\textnormal{N}}}
\def\ermO{{\textnormal{O}}}
\def\ermP{{\textnormal{P}}}
\def\ermQ{{\textnormal{Q}}}
\def\ermR{{\textnormal{R}}}
\def\ermS{{\textnormal{S}}}
\def\ermT{{\textnormal{T}}}
\def\ermU{{\textnormal{U}}}
\def\ermV{{\textnormal{V}}}
\def\ermW{{\textnormal{W}}}
\def\ermX{{\textnormal{X}}}
\def\ermY{{\textnormal{Y}}}
\def\ermZ{{\textnormal{Z}}}

% Vectors
\def\vzero{{\bm{0}}}
\def\vone{{\bm{1}}}
\def\vmu{{\bm{\mu}}}
\def\vtheta{{\bm{\theta}}}
\def\va{{\bm{a}}}
\def\vb{{\bm{b}}}
\def\vc{{\bm{c}}}
\def\vd{{\bm{d}}}
\def\ve{{\bm{e}}}
\def\vf{{\bm{f}}}
\def\vg{{\bm{g}}}
\def\vh{{\bm{h}}}
\def\vi{{\bm{i}}}
\def\vj{{\bm{j}}}
\def\vk{{\bm{k}}}
\def\vl{{\bm{l}}}
\def\vm{{\bm{m}}}
\def\vn{{\bm{n}}}
\def\vo{{\bm{o}}}
\def\vp{{\bm{p}}}
\def\vq{{\bm{q}}}
\def\vr{{\bm{r}}}
\def\vs{{\bm{s}}}
\def\vt{{\bm{t}}}
\def\vu{{\bm{u}}}
\def\vv{{\bm{v}}}
\def\vw{{\bm{w}}}
\def\vx{{\bm{x}}}
\def\vy{{\bm{y}}}
\def\vz{{\bm{z}}}

% Elements of vectors
\def\evalpha{{\alpha}}
\def\evbeta{{\beta}}
\def\evepsilon{{\epsilon}}
\def\evlambda{{\lambda}}
\def\evomega{{\omega}}
\def\evmu{{\mu}}
\def\evpsi{{\psi}}
\def\evsigma{{\sigma}}
\def\evtheta{{\theta}}
\def\eva{{a}}
\def\evb{{b}}
\def\evc{{c}}
\def\evd{{d}}
\def\eve{{e}}
\def\evf{{f}}
\def\evg{{g}}
\def\evh{{h}}
\def\evi{{i}}
\def\evj{{j}}
\def\evk{{k}}
\def\evl{{l}}
\def\evm{{m}}
\def\evn{{n}}
\def\evo{{o}}
\def\evp{{p}}
\def\evq{{q}}
\def\evr{{r}}
\def\evs{{s}}
\def\evt{{t}}
\def\evu{{u}}
\def\evv{{v}}
\def\evw{{w}}
\def\evx{{x}}
\def\evy{{y}}
\def\evz{{z}}

% Matrix
\def\mA{{\bm{A}}}
\def\mB{{\bm{B}}}
\def\mC{{\bm{C}}}
\def\mD{{\bm{D}}}
\def\mE{{\bm{E}}}
\def\mF{{\bm{F}}}
\def\mG{{\bm{G}}}
\def\mH{{\bm{H}}}
\def\mI{{\bm{I}}}
\def\mJ{{\bm{J}}}
\def\mK{{\bm{K}}}
\def\mL{{\bm{L}}}
\def\mM{{\bm{M}}}
\def\mN{{\bm{N}}}
\def\mO{{\bm{O}}}
\def\mP{{\bm{P}}}
\def\mQ{{\bm{Q}}}
\def\mR{{\bm{R}}}
\def\mS{{\bm{S}}}
\def\mT{{\bm{T}}}
\def\mU{{\bm{U}}}
\def\mV{{\bm{V}}}
\def\mW{{\bm{W}}}
\def\mX{{\bm{X}}}
\def\mY{{\bm{Y}}}
\def\mZ{{\bm{Z}}}
\def\mBeta{{\bm{\beta}}}
\def\mPhi{{\bm{\Phi}}}
\def\mLambda{{\bm{\Lambda}}}
\def\mSigma{{\bm{\Sigma}}}

% Tensor
\DeclareMathAlphabet{\mathsfit}{\encodingdefault}{\sfdefault}{m}{sl}
\SetMathAlphabet{\mathsfit}{bold}{\encodingdefault}{\sfdefault}{bx}{n}
\newcommand{\tens}[1]{\bm{\mathsfit{#1}}}
\def\tA{{\tens{A}}}
\def\tB{{\tens{B}}}
\def\tC{{\tens{C}}}
\def\tD{{\tens{D}}}
\def\tE{{\tens{E}}}
\def\tF{{\tens{F}}}
\def\tG{{\tens{G}}}
\def\tH{{\tens{H}}}
\def\tI{{\tens{I}}}
\def\tJ{{\tens{J}}}
\def\tK{{\tens{K}}}
\def\tL{{\tens{L}}}
\def\tM{{\tens{M}}}
\def\tN{{\tens{N}}}
\def\tO{{\tens{O}}}
\def\tP{{\tens{P}}}
\def\tQ{{\tens{Q}}}
\def\tR{{\tens{R}}}
\def\tS{{\tens{S}}}
\def\tT{{\tens{T}}}
\def\tU{{\tens{U}}}
\def\tV{{\tens{V}}}
\def\tW{{\tens{W}}}
\def\tX{{\tens{X}}}
\def\tY{{\tens{Y}}}
\def\tZ{{\tens{Z}}}


% Graph
\def\gA{{\mathcal{A}}}
\def\gB{{\mathcal{B}}}
\def\gC{{\mathcal{C}}}
\def\gD{{\mathcal{D}}}
\def\gE{{\mathcal{E}}}
\def\gF{{\mathcal{F}}}
\def\gG{{\mathcal{G}}}
\def\gH{{\mathcal{H}}}
\def\gI{{\mathcal{I}}}
\def\gJ{{\mathcal{J}}}
\def\gK{{\mathcal{K}}}
\def\gL{{\mathcal{L}}}
\def\gM{{\mathcal{M}}}
\def\gN{{\mathcal{N}}}
\def\gO{{\mathcal{O}}}
\def\gP{{\mathcal{P}}}
\def\gQ{{\mathcal{Q}}}
\def\gR{{\mathcal{R}}}
\def\gS{{\mathcal{S}}}
\def\gT{{\mathcal{T}}}
\def\gU{{\mathcal{U}}}
\def\gV{{\mathcal{V}}}
\def\gW{{\mathcal{W}}}
\def\gX{{\mathcal{X}}}
\def\gY{{\mathcal{Y}}}
\def\gZ{{\mathcal{Z}}}

% Sets
\def\sA{{\mathbb{A}}}
\def\sB{{\mathbb{B}}}
\def\sC{{\mathbb{C}}}
\def\sD{{\mathbb{D}}}
% Don't use a set called E, because this would be the same as our symbol
% for expectation.
\def\sF{{\mathbb{F}}}
\def\sG{{\mathbb{G}}}
\def\sH{{\mathbb{H}}}
\def\sI{{\mathbb{I}}}
\def\sJ{{\mathbb{J}}}
\def\sK{{\mathbb{K}}}
\def\sL{{\mathbb{L}}}
\def\sM{{\mathbb{M}}}
\def\sN{{\mathbb{N}}}
\def\sO{{\mathbb{O}}}
\def\sP{{\mathbb{P}}}
\def\sQ{{\mathbb{Q}}}
\def\sR{{\mathbb{R}}}
\def\sS{{\mathbb{S}}}
\def\sT{{\mathbb{T}}}
\def\sU{{\mathbb{U}}}
\def\sV{{\mathbb{V}}}
\def\sW{{\mathbb{W}}}
\def\sX{{\mathbb{X}}}
\def\sY{{\mathbb{Y}}}
\def\sZ{{\mathbb{Z}}}

% Entries of a matrix
\def\emLambda{{\Lambda}}
\def\emA{{A}}
\def\emB{{B}}
\def\emC{{C}}
\def\emD{{D}}
\def\emE{{E}}
\def\emF{{F}}
\def\emG{{G}}
\def\emH{{H}}
\def\emI{{I}}
\def\emJ{{J}}
\def\emK{{K}}
\def\emL{{L}}
\def\emM{{M}}
\def\emN{{N}}
\def\emO{{O}}
\def\emP{{P}}
\def\emQ{{Q}}
\def\emR{{R}}
\def\emS{{S}}
\def\emT{{T}}
\def\emU{{U}}
\def\emV{{V}}
\def\emW{{W}}
\def\emX{{X}}
\def\emY{{Y}}
\def\emZ{{Z}}
\def\emSigma{{\Sigma}}

% entries of a tensor
% Same font as tensor, without \bm wrapper
\newcommand{\etens}[1]{\mathsfit{#1}}
\def\etLambda{{\etens{\Lambda}}}
\def\etA{{\etens{A}}}
\def\etB{{\etens{B}}}
\def\etC{{\etens{C}}}
\def\etD{{\etens{D}}}
\def\etE{{\etens{E}}}
\def\etF{{\etens{F}}}
\def\etG{{\etens{G}}}
\def\etH{{\etens{H}}}
\def\etI{{\etens{I}}}
\def\etJ{{\etens{J}}}
\def\etK{{\etens{K}}}
\def\etL{{\etens{L}}}
\def\etM{{\etens{M}}}
\def\etN{{\etens{N}}}
\def\etO{{\etens{O}}}
\def\etP{{\etens{P}}}
\def\etQ{{\etens{Q}}}
\def\etR{{\etens{R}}}
\def\etS{{\etens{S}}}
\def\etT{{\etens{T}}}
\def\etU{{\etens{U}}}
\def\etV{{\etens{V}}}
\def\etW{{\etens{W}}}
\def\etX{{\etens{X}}}
\def\etY{{\etens{Y}}}
\def\etZ{{\etens{Z}}}

% The true underlying data generating distribution
\newcommand{\pdata}{p_{\rm{data}}}
% The empirical distribution defined by the training set
\newcommand{\ptrain}{\hat{p}_{\rm{data}}}
\newcommand{\Ptrain}{\hat{P}_{\rm{data}}}
% The model distribution
\newcommand{\pmodel}{p_{\rm{model}}}
\newcommand{\Pmodel}{P_{\rm{model}}}
\newcommand{\ptildemodel}{\tilde{p}_{\rm{model}}}
% Stochastic autoencoder distributions
\newcommand{\pencode}{p_{\rm{encoder}}}
\newcommand{\pdecode}{p_{\rm{decoder}}}
\newcommand{\precons}{p_{\rm{reconstruct}}}

\newcommand{\laplace}{\mathrm{Laplace}} % Laplace distribution

\newcommand{\E}{\mathbb{E}}
\newcommand{\Ls}{\mathcal{L}}
\newcommand{\R}{\mathbb{R}}
\newcommand{\emp}{\tilde{p}}
\newcommand{\lr}{\alpha}
\newcommand{\reg}{\lambda}
\newcommand{\rect}{\mathrm{rectifier}}
\newcommand{\softmax}{\mathrm{softmax}}
\newcommand{\sigmoid}{\sigma}
\newcommand{\softplus}{\zeta}
\newcommand{\KL}{D_{\mathrm{KL}}}
\newcommand{\Var}{\mathrm{Var}}
\newcommand{\standarderror}{\mathrm{SE}}
\newcommand{\Cov}{\mathrm{Cov}}
% Wolfram Mathworld says $L^2$ is for function spaces and $\ell^2$ is for vectors
% But then they seem to use $L^2$ for vectors throughout the site, and so does
% wikipedia.
\newcommand{\normlzero}{L^0}
\newcommand{\normlone}{L^1}
\newcommand{\normltwo}{L^2}
\newcommand{\normlp}{L^p}
\newcommand{\normmax}{L^\infty}

\newcommand{\parents}{Pa} % See usage in notation.tex. Chosen to match Daphne's book.

\DeclareMathOperator*{\argmax}{arg\,max}
\DeclareMathOperator*{\argmin}{arg\,min}

\DeclareMathOperator{\sign}{sign}
\DeclareMathOperator{\Tr}{Tr}
\let\ab\allowbreak


\usepackage{url}
\usepackage{graphicx}
\usepackage{booktabs}
\usepackage{array}
\usepackage{longtable}
\usepackage{float}
\usepackage[section]{placeins}
\usepackage[hidelinks]{hyperref}
\newcolumntype{P}[1]{>{\raggedright\arraybackslash}p{#1}}
\setlength{\emergencystretch}{2em}

\title{Supplementary Material for\\
Beyond Correctness: Evaluating Semantic Knowledge in Cross-Table Transfer}
\author{Anonymous Authors}

\begin{document}
\maketitle
\def\ICLRFinishAppendix{%
\bibliography{references_seed}%
\bibliographystyle{iclr2027_conference}%
\end{document}%
}
\fi

\appendix
\raggedbottom
\renewcommand{\thefigure}{A\arabic{figure}}
\renewcommand{\thetable}{A\arabic{table}}
\renewcommand{\theequation}{A\arabic{equation}}
\ifdefined\theHfigure
\renewcommand{\theHfigure}{appendix.\arabic{figure}}
\renewcommand{\theHtable}{appendix.\arabic{table}}
\renewcommand{\theHequation}{appendix.\arabic{equation}}
\fi
\setcounter{figure}{0}
\setcounter{table}{0}
\setcounter{equation}{0}


\section{Supplementary Experimental and Construction Details}
\label{sec:app-scope}

This appendix provides additional experimental details, extended results, and robustness analyses for the experiments in the main paper.
Unless stated otherwise, the reported results are specific to the evaluated datasets, learners, interventions, and references.
The accompanying artifact contains code, frozen semantic records and reference libraries, prompt and generation records for the LLM-based construction step, analysis manifests, environment information, and mappings from reported tables and figures to their generating scripts and outputs.

\subsection{Common experimental and inference settings}
\label{sec:app-experimental-details}

\paragraph{Comparisons and notation.}
Content sensitivity compares the intended and altered conditions, whereas predictive utility compares the intended condition with a suitable reference that removes the tested semantic knowledge while preserving the other predictive information.
Target-only models are reported separately when relevant to show performance without source transfer; they are not used as references for individual semantic components.

We use the following decomposition:

$$
Q_{\mathrm{intended}} - Q_{\mathrm{altered}}
=
\left(
Q_{\mathrm{intended}} - Q_{\mathrm{reference}}
\right)
+
\left(
Q_{\mathrm{reference}} - Q_{\mathrm{altered}}
\right).
$$

The first term on the right is predictive utility.
The second term is reported only to show how much of the intended--altered difference comes from the reference--altered difference; we do not treat it as a separate estimand.


\paragraph{Scoring and reporting conventions.}
Intended--altered and intended--reference differences are reported so that positive values favor the intended condition.
The controlled benchmark uses
\[
Q=-\mathrm{MSE}/\mathrm{Var}(y^{\mathrm{qry}}).
\]

Absolute performance is reported as
\[
\mathrm{nMSE}
=
\mathrm{MSE}/\mathrm{Var}(y^{\mathrm{qry}})
=
-Q.
\]
Real-data experiments use their original AUROC, nRMSE, or log-RMSE metrics, so numerical effect sizes are compared only within the same metric.

Whenever nRMSE is reported, it is defined separately for each endpoint,
transfer direction, and split using the held-out target query set \(Q\):
\[
\operatorname{RMSE}_Q
=
\sqrt{
\frac{1}{|Q|}
\sum_{i\in Q}
(y_i-\hat y_i)^2
},
\qquad
s_Q
=
\sqrt{
\frac{1}{|Q|}
\sum_{i\in Q}
(y_i-\bar y_Q)^2
},
\]
and
\[
\operatorname{nRMSE}_Q
=
\frac{\operatorname{RMSE}_Q}{s_Q}.
\]
Here \(s_Q\) is the population standard deviation of the gold labels in the
same held-out target query set (\(\mathrm{ddof}=0\)).
Query labels are used only for evaluation and normalization, not for model fitting.
Lower nRMSE is better; under the common higher-is-better convention,
\(Q=-\mathrm{nRMSE}\), so intended--reference utility is equivalently
\(\mathrm{nRMSE}_{\mathrm{reference}}
-\mathrm{nRMSE}_{\mathrm{intended}}\).

This definition is distinct from the controlled benchmark's primary
normalized error,
\(\mathrm{nMSE}=\mathrm{MSE}/\mathrm{Var}(y^{\mathrm{qry}})\).

\paragraph{Common experimental settings.}
Labeled target rows used for training are separate from the query rows, and query labels are used only for evaluation.
Within each comparison, conditions use the same data split, source sample, labeled-target sample, and random seed whenever applicable.

\begin{table}[t]
\caption{Common experimental settings across the real-data and controlled experiments.}
\label{tab:app-common-settings}
\centering
\footnotesize
\setlength{\tabcolsep}{2pt}

\begin{tabular}{@{}P{0.13\linewidth}P{0.18\linewidth}P{0.12\linewidth}P{0.225\linewidth}P{0.14\linewidth}P{0.115\linewidth}@{}}

\toprule
Setting & Target setup & Model & Model settings & Runs & Inference unit \\
\midrule

NHANES--\newline KNHANES &
Source rows plus \(n_{\mathrm{target}}=256\) labeled target rows; disjoint target query &
XGB &
300 trees, depth 6, learning rate .05; source and target total weights balanced &
10 seeds (50--59) per endpoint &
Endpoint \\
\midrule

NHANES/\newline KNHANES--HRS &
4,096 capped source rows; 1/5/10\% labeled target rows &
XGB &
300 trees, depth 5, learning rate .05; target weight 10 &
10 seeds (40--49) per labeled-target fraction &
Endpoint \\
\midrule

CAMELS-US--\newline CAMELS-GB &
1/3/5 labeled target basins; fixed disjoint query basins &
XGB &
300 trees, depth 6, learning rate .05; log-target regression; target weight 10 &
20 seeds (20--39) per labeled-target size &
Basin \\
\midrule

Controlled &
2,048 source rows; \(n_{\mathrm{target}}\in\{32,512\}\); 2,000 query rows &
XGB, HistGB, TabPFN, MLP &
Primary conditions use unit weights; XGB/HistGB: 300 depth-6 iterations; TabPFN 7.0.0 defaults; MLP: one 32-unit ReLU layer &
20 realizations \(\times\) 3 fixed families &
Realization within family \\

\bottomrule
\end{tabular}
\end{table}
Model-specific procedures for the retrospective TabLLM experiments and the
TransTab/CARTE learner-coverage extensions are reported in
Sections~\ref{sec:app-tabllm-reproduction} and
\ref{sec:app-learner-coverage}, respectively.
For conditions that use source data, source and labeled-target rows are concatenated and fit with the same model configuration.
Target-only models use only the labeled target rows and are reported separately when relevant.

\paragraph{Pairing and intervals.}
Conditions within the same comparison use the same data split and random seed, and query performance is never used to select hyperparameters.
Unless otherwise stated, bootstrap intervals use 10,000 resamples. The paired reference-variation analysis in Section~\ref{sec:app-tabllm-reference-audit} uses 5,000 resamples, and the fixed-mapping lexical-reference analysis in Section C.3 uses 2,000 paired resamples.
For the survey experiments, we resample endpoints and their paired seeds; for hydrology, we resample basins; and for the controlled experiments, we resample the 20 realizations within each outcome family.
When an analysis uses a different resampling procedure or decision rule, we state it in that section.
Alternative interval estimates are reported only as robustness checks and do not replace the primary result.
The reported intervals apply to each contrast separately; we do not adjust them jointly across all analyses.
Accordingly, resolved and unresolved refer to the stated contrast rather than to a joint discovery claim.
When many references are compared at once, we use simultaneous intervals for the sign-change analysis.

\subsection{Analysis provenance}
\label{sec:app-registry}

We distinguish analyses by when their design was fixed relative to the results they evaluate.
\emph{Prespecified} analyses were fixed before the corresponding comparison;
\emph{frozen extensions} were added after earlier analyses but fixed before their new execution;
\emph{post-result} analyses were developed after relevant earlier results were known; and
\emph{retrospective} analyses operate on previously published or released results.
Replication and exploratory labels describe analysis scope rather than timing.
The complete analysis ledger, including construction dates and decision records, is retained in the accompanying artifact.

\begin{table}[h]
\caption{
Provenance information for analyses whose timing affects interpretation.
}
\label{tab:app-analysis-registry}
\centering
\footnotesize
\setlength{\tabcolsep}{3pt}

\begin{tabular}{P{0.30\linewidth}P{0.25\linewidth}P{0.37\linewidth}}
\toprule
Analysis & Timing / status & Interpretation \\
\midrule

Measurement endpoint extension
& Frozen extension
& The broader ten-endpoint result is descriptive. \\

Examination width-matched reference
& Post-result construction
& Enables a new predictive-utility comparison; it does not retroactively make the original ablation a suitable reference comparison. \\

Examination target-support ladder
& Direction calibrated earlier; six-endpoint panel selected after the 13-endpoint expansion
& Six-endpoint trajectory is treated as replication; the full 13-endpoint trajectory is exploratory. \\

Relation capacity analysis
& Frozen extension
& Evaluates learner-capacity dependence of relation utility. \\

Fixed-width noisy-proxy diagnostic
& Post-result diagnostic; protocol fixed before execution
& Tests reconstructibility while preserving frame width and variable identity; the within-dose association is post-result descriptive. \\

Real-data reconstructibility probes
& Frozen extension
& The reconstruction-manipulation criterion was not met in the evaluated real-data settings. \\

TabLLM reproduction
& Retrospective compatibility reproduction
& Reanalyzes released predictions and code under a format-matched reference. \\

TabLLM reference-space audit
& Frozen extension
& The primary anonymous reference predates the audit; additional references evaluate sensitivity rather than reselecting the primary reference. \\

\bottomrule
\end{tabular}
\end{table}

\subsection{Controlled benchmark construction}

\paragraph{Data and outcomes.}
The controlled benchmark uses an eight-feature NHANES panel while preserving the observed feature values and missingness patterns.
Source participants are drawn from the 2001--2014 cycles and target participants from the 2015--2018 and 2023 cycles, with no participant overlap between the two groups.

Missing entries are filled only for synthetic-outcome construction using draws based on the observed source distribution; the learner continues to receive the original inputs with their original missingness.

For source-standardized mechanism terms $\widetilde h_t(z_i)$, outcomes are generated as
\begin{equation}
y_i = \eta_i + \epsilon_i,
\qquad
\eta_i = \sum_{t\in\mathcal T} w_t s_t \widetilde h_t(z_i),
\qquad
\epsilon_i \sim
\mathcal N\!\left(0,\operatorname{SD}_s(\eta)^2\right),
\end{equation}
where $w_t\sim U(.5,1)$ and $s_t\in\{-1,+1\}$ are fixed within each realization.
This noise construction gives a source signal-to-noise ratio of one.

We use three outcome families.
The additive family selects five single-feature terms.
The pairwise family selects three disjoint feature pairs over six features, and the sparse family selects two pairs from all 28 possible pairs.
Pairwise terms are either feature differences or log-differences,
with standardized inputs clipped to $[-5,5]$ before the logarithmic transformation.

Each realization fixes the outcome mechanism, schema rendering, source and target rows, labeled-target sample, and noise before learner comparison.
We generate 20 realizations for each outcome family.

\paragraph{Schema differences and altered conditions.}
Four features are represented as ordinal variables and four as continuous variables.
Across the source and target schemas, ordinal features can differ in coding direction, code base, and missing-value sentinel, while continuous features can differ in scale, offset, and sentinel.
The target schema is kept fixed across all conditions.

For measurement knowledge, the altered condition changes the supplied coding information by permuting ordinal values, assigning continuous transformations to the wrong features, and changing sentinel information.
For correspondence knowledge, the altered condition reassigns features within the ordinal and continuous groups so that no feature remains matched to itself.
For relation knowledge, the altered condition replaces the intended relations with different relations while keeping the number and type of relations comparable.
Reversed and randomly replaced relations are evaluated separately.

\paragraph{Mismatch dose.}
For the measurement experiments, we also vary how many source features retain their dataset-specific schema.
Each realization uses a fixed ordering of the four ordinal and four continuous features.
At dose zero, all source features use the canonical representation.
As the dose increases, more source features use their dataset-specific representation, while the labeled target and query data remain unchanged.

\subsection{Evaluated feature and endpoint panels}
\label{sec:app-evaluated-panels}

\paragraph{Controlled feature panel.}
The controlled benchmark uses eight NHANES variables:
age, total cholesterol, HbA1c, creatinine, hemoglobin,
red blood cell count, white blood cell count, and waist circumference.
In each realization, four of the eight variables are randomly selected and
rendered as five-level ordinal variables using source quantiles, while the
remaining four are treated as continuous variables.
Thus, the ordinal subset is not tied to a fixed set of variable identities
and varies across realizations.

\paragraph{Survey-label measurement panel.}
The measurement experiments treat documented questionnaire codes as binary
prediction endpoints.
The focused panel contains diabetes history, hypertension history, and current
smoking status; the broader extension adds kidney disease, stroke, arthritis,
asthma, cancer, heart attack, and angina history.
The endpoint's own item is excluded from the predictor set.

\begin{table}[h]
\caption{
Survey-label endpoints used in the NHANES--KNHANES measurement experiments.
The first three rows form the focused panel; all ten rows form the extended panel.
}
\label{tab:app-measurement-endpoints}
\centering
\footnotesize
\setlength{\tabcolsep}{3pt}

\begin{tabular}{P{0.35\linewidth}P{0.22\linewidth}P{0.32\linewidth}}
\toprule
Endpoint & NHANES item & KNHANES item \\
\midrule
Diabetes history
& \texttt{DIQ010}
& \texttt{ALL\_\_de1\_dg} \\

Hypertension history
& \texttt{BPQ020}
& \texttt{ALL\_\_di1\_dg} \\

Current smoking status
& \texttt{SMQ040}
& \texttt{ALL\_\_bs3\_1} \\

Kidney disease history
& \texttt{KIQ022}
& \texttt{ALL\_\_dn1\_dg} \\

Stroke history
& \texttt{MCQ160F}
& \texttt{ALL\_\_di3\_dg} \\

Arthritis history
& \texttt{MCQ160A}
& \texttt{ALL\_\_dm1\_dg} \\

Asthma history
& \texttt{MCQ010}
& \texttt{ALL\_\_dj4\_dg} \\

Cancer history
& \texttt{MCQ220}
& \texttt{ALL\_\_dc1\_dg} \\

Heart attack history
& \texttt{MCQ160E}
& \texttt{ALL\_\_di5\_dg} \\

Angina history
& \texttt{MCQ160D}
& \texttt{ALL\_\_di6\_dg} \\

\bottomrule
\end{tabular}
\end{table}

\paragraph{Examination panel.}
The examination analysis contains 13 eligible endpoints:
glucose, waist circumference, triglycerides, systolic blood pressure,
diastolic blood pressure, total cholesterol, alanine aminotransferase (ALT),
blood urea nitrogen (BUN), creatinine, hemoglobin, hematocrit,
red blood cell count, and white blood cell count.
The original six-endpoint panel consists of glucose, waist circumference,
triglycerides, systolic and diastolic blood pressure, and total cholesterol.
The remaining seven endpoints were added in the broader endpoint extension.
Pulse rate was excluded by the frozen requirement of at least 1,024 eligible
target-domain rows.

The examination predictor registry contains 20 candidate slots.
The 11 base slots are age, creatinine, red blood cell count,
white blood cell count, hemoglobin, hematocrit, sex, education level,
diabetes history, hypertension history, and current smoking status.
The nine correspondence-member slots are height, weight, waist circumference,
systolic blood pressure, diastolic blood pressure, glucose, HbA1c,
total cholesterol, and triglycerides.
For each task, the target endpoint itself is removed from the predictor set.

\subsection{Additional details on suitable references}
\label{sec:app-reference-checks}

The main paper defines a suitable reference using three criteria: Removal, Preservation, and Interface.
Table~\ref{tab:app-reference-portability} shows how these criteria apply to the comparisons used in the measurement, correspondence, and relation experiments.
A comparison that fails one or more criteria may still be useful as an altered condition or descriptive ablation, but it is not used to estimate predictive utility.

\begin{table}[h]
\caption{
Application of the reference criteria across the measurement, correspondence, and relation experiments.
All listed comparisons satisfy the Interface criterion.
}
\label{tab:app-reference-portability}
\centering
\small
\setlength{\tabcolsep}{4pt}

\begin{tabular}{P{0.38\linewidth}ccP{0.34\linewidth}}
\toprule
Comparison & Removal & Preservation & Interpretation \\
\midrule

Measurement placebo
& $\times$
& \checkmark
& Altered condition \\
\midrule

Controlled correspondence derangement
& $\times$
& \checkmark
& Content sensitivity only \\
\midrule

Survey no-added-columns
& \checkmark
& $\times$
& Descriptive ablation \\
\midrule

Examination domain-local
& \checkmark
& \checkmark
& Suitable reference \\
\midrule

Relation flipped / matched-random
& $\times$
& \checkmark
& Altered conditions \\
\midrule

Relation unconstrained reference
& \checkmark
& \checkmark
& Suitable reference \\

\bottomrule
\end{tabular}
\end{table}

The survey no-added-columns comparison satisfies Removal but not Preservation because it also removes other predictive information.
The controlled correspondence derangement satisfies Preservation but not Removal because it changes the correspondence rather than removing it.
We therefore treat the former as a descriptive ablation and the latter as a content-sensitivity comparison.


\subsection{Semantic-knowledge construction and audits}
\label{sec:app-extraction-documentation}

\paragraph{Primary LLM-based candidate generation.}

For the NHANES--KNHANES documentation-based proposal run, we used GPT-5.6 Sol through Codex CLI 0.144.5 with reasoning effort set to \texttt{high}.
The system instruction and user template are stored separately for provenance, but during execution they were concatenated and sent as a single CLI input rather than as separate message roles.

The model receives schema metadata and selected documentation spans, but not person-level rows, target labels, predictions, or model-performance results.
The prompt also contains the fixed output JSON schema and the executable relation DSL.
Concept candidates may group multiple documented variables across the source and target schemas and are not restricted to one-to-one matches.
Relation candidates must use an operation from the fixed DSL and cite documentation for both source and target bindings.

After generation, fixed mechanical checks are applied to the output.
These checks cover JSON and schema validity, eligible columns, evidence references, prohibited leakage roles, operation and argument requirements, unit compatibility, duplicate candidates, and fixed candidate budgets.
LLM confidence, downstream predictive performance, and human semantic judgments are not used to accept candidates.

The stored generation configuration lists greedy decoding, temperature 0, and seed 0.
These settings, however, do not appear as explicit options in the recorded Codex CLI command, so we do not claim deterministic decoding.
Instead, we preserve the system instruction, user template, rendered prompt, input and output schemas, validation rules, raw response, execution manifest, and cryptographic hashes for the executed run.
The experiments therefore use the stored candidate set from this recorded run; reproducibility does not require a new LLM call to generate the same candidates.

\paragraph{Seven-model measurement-code extraction audit.}

Separately from the candidate-generation procedure above, we evaluate measurement-code extraction from documentation across seven language models.
This audit uses its own prompt and validation procedure and is not a rerun of the GPT-5.6 Sol candidate-generation procedure.
Each model is run once with the fixed audit prompt, without external tools or retrieval, and the raw output is preserved.

Across seven evaluated models, five produced complete records in the required format, and four also produced valid-code sets that covered the target codes needed for evaluation.
The purpose of this analysis is to check whether measurement-code extraction can be carried out consistently from documentation alone, rather than to evaluate predictive utility.

\paragraph{Correspondence construction audit.}
\label{sec:app-correspondence-construction-audit}

Separately, we test correspondence construction in a larger candidate space containing 1,315 target variables.
Each of nine known source--target matches is evaluated against this full candidate set. The artifact includes a metadata-only export of all 1,315 column names; its source-file hash and sorted-name pool hash match the frozen execution record. The CSV was exported during artifact finalization, not represented as a historical CSV.

Using feature names alone recovers 2 of the 9 matches, while adding units recovers 4 of 9.
Using names, units, and codebook descriptions recovers all 9 matches.
When the true match is removed from the candidate set, a single score threshold retains 8 of the 9 true matches in the standard setting without assigning a false match in any of the nine corresponding no-match cases.

For comparison, the closed 9-by-9 matching task recovers all 9 matches.
This indicates that the smaller closed task provides an easier correspondence setting than matching against the full target schema.

\paragraph{Evidence for the relation set.}
\label{sec:app-relation-evidence}

The relations used in the real-data experiments were specified before predictive evaluation.
We subsequently reviewed prior literature to document the evidence for each relation and how closely that evidence matches the variables and prediction setting used here.
This review was used only to characterize the selected relations; it did not change the relation set.

\begin{table}[h]
\caption{External evidence for the relations used in the real-data experiments. The cited evidence supports the stated direction to different degrees and does not imply that each relation is a universal constraint.}
\label{tab:app-relation-evidence}
\centering
\footnotesize
\setlength{\tabcolsep}{3pt}

\begin{tabular}{P{0.20\linewidth}P{0.07\linewidth}P{0.27\linewidth}P{0.38\linewidth}}
\toprule
Relation & Dir. & Supporting evidence & Scope of the evidence \\
\midrule

Age $\rightarrow$ SBP
& $+$
& \citep{franklin1997hemodynamic,mitchell2010hemodynamic}
& Population trends do not imply the same association under every conditioning set. \\
\midrule

SBP $\leftrightarrow$ DBP
& $+$
& \citep{zhu2022pleiotropy,franklin1997hemodynamic}
& The population association is positive, although DBP changes differently with age. \\
\midrule

Fasting glucose $\leftrightarrow$ HbA1c
& $+$
& \citep{nathan2008translating}
& The cited relationship uses average rather than fasting glucose. \\
\midrule

Total cholesterol $\leftrightarrow$ triglycerides
& $+$
& \citep{friedewald1972estimation,wolska2026newmethods}
& Their joint use in LDL estimation is not direct evidence of a general population correlation. \\
\midrule

Weight $\rightarrow$ waist circumference
& $+$
& \citep{flegal2009comparisons,ross2020waist}
& Much of the cited evidence concerns BMI rather than weight alone. \\
\midrule

BUN $\leftrightarrow$ creatinine
& $+$
& \citep{hosten1990bun,morgan1977creatinine}
& Both reflect renal function, but their association is not uniformly strong across populations. \\
\midrule

Hemoglobin $\leftrightarrow$ hematocrit/RBC
& $+$
& \citep{doig2017methodical,guevara2023hemoglobin}
& The relationship depends on red-cell indices and population characteristics. \\
\midrule

Precipitation $\rightarrow$ discharge
& $+$
& \citep{zhang2001response,addor2017camels,coxon2020camelsgb,andreassian2025timeshift}
& Much of the cited evidence concerns annual or long-term behavior, whereas prediction uses daily data. \\
\midrule

PET $\rightarrow$ discharge
& $-$
& \citep{zhang2001response,addor2017camels,coxon2020camelsgb,andreassian2025timeshift}
& The long-term relationship is weak and not uniformly negative, and is applied here at a daily scale. \\

\bottomrule
\end{tabular}
\end{table}

\section{Additional Evidence on Content Sensitivity and Predictive Utility}

This section provides additional results comparing content sensitivity with predictive utility across the evaluated settings.


\subsection{TabLLM feature-name ablation}
\label{sec:app-tabllm-reproduction}

We revisit the published TabLLM feature-name ablation as an example in which the intended--altered difference can be compared with performance relative to a reference.

In the published zero-shot results, the mean intended--altered difference across the nine datasets is \(+.107\) AUROC.
The corresponding difference between the original List Template and the published values-only condition is \(+.076\).
The values-only condition, however, removes both feature names and the surrounding template structure, so we treat it as a published comparator rather than a suitable reference for feature-name utility.

The dataset-level results already show that the two differences need not agree.
For Credit-g, the intended--altered difference is \(+.09\), whereas the original names perform \(-.13\) below the values-only condition.
Thus, the large intended--altered difference in this case does not by itself show that feature names improve prediction.

To evaluate predictive utility while preserving the remaining input format, we reproduce the released TabLLM evaluation using a reference with fixed anonymous feature identifiers.
This reference keeps the List Template, field order, value formatting, and feature identity unchanged while replacing the feature names with identifiers such as \texttt{feature\_001}.
The identifiers keep the fields distinct without supplying their original feature meanings; differences among identifier strings are examined separately in Section~\ref{sec:app-tabllm-reference-audit}.

\paragraph{Compatibility reproduction.}
We treat this experiment as a released-code compatibility reproduction rather
than a bit-exact reproduction of the historical software stack.
The released 11B execution configuration and IA3 checkpoint resolve to
\texttt{bigscience/T0}, although one description in the original paper refers
to T0pp; we follow the executable released configuration and checkpoint.

Because the released repository contains only part of the corresponding
T-Few implementation and the historical software stack does not execute
directly in our current environment, we use the model-modification code from
the pinned corresponding T-Few revision together with documented dependency
and bfloat16 compatibility fixes.
These changes preserve the released model initialization, serialization,
scoring procedure, and semantic comparison.

Under this compatibility reproduction, 21 of the 27 published-condition
dataset means are within .015 AUROC of the paper's printed values, with a
maximum absolute deviation of .0347.
We therefore do not claim a bit-exact numerical reproduction.
All semantic contrasts reported here are computed pairwise within this common
reproduction.

\paragraph{Adaptation protocol.}
We evaluate nine datasets---Bank, Blood, California Housing, Car, Credit-g,
Diabetes, Heart, Income, and Jungle---at
\(n_{\mathrm{adapt}}\in\{0,4,32,512\}\)
using the released split seeds
\(42, 1024, 0, 1, 32\).
The \(n_{\mathrm{adapt}}=0\) condition reuses the frozen zero-shot compatibility
reproduction without task-specific fitting.

For \(n_{\mathrm{adapt}}>0\), we initialize the released
\texttt{bigscience/T0} 11B model from the public IA3 checkpoint and update only
the 192 IA3 \texttt{lora\_b} tensors
(1,081,344 trainable parameters), while keeping the T0 backbone fixed.
We use the released T-Few objective consisting of the language-model,
multiple-choice, and unlikely losses.
Optimization uses Adafactor with learning rate \(0.003\), zero weight decay,
linear decay with 6\% warmup, gradient clipping at 1, batch size 4,
and 30 epochs.
This gives 30, 240, and 3,840 update steps for
\(n_{\mathrm{adapt}}=4,32,512\), respectively.
The maximum input length is 1,024 tokens.

For each seed, we construct the released-style 80/20 train--test split.
Adaptation examples are sampled class-balanced with replacement from the
remaining 80\%.
Paired semantic conditions use identical split membership, adaptation-row
membership, replacement multiplicities, row order, and loader seed.
Prediction uses normalized class-verbalizer sequence likelihoods;
binary tasks are evaluated by class-1 AUROC and Car by macro one-vs-rest AUROC.

The historical PyTorch-Lightning orchestration is replaced by a standalone
PyTorch compatibility port.
Sampling, tokenization, objective, optimizer, schedule, update count,
checkpoint initialization, and scoring are held to the released T-Few
contract and were checked by a frozen zero-shot conformance test before the
adaptation experiments.


For TabLLM, intervals use 10,000 paired multiplier resamples over the shared test examples, with dataset resampling for the macro estimate.
Zero-shot evaluations on the same test examples are treated as paired observations rather than independent runs.

\begin{table}[h]
\caption{
Released-code TabLLM reproduction using the anonymous-name reference.
Values are AUROC differences with 95\% paired intervals.
}
\label{tab:app-tabllm-reproduction}
\centering
\footnotesize
\setlength{\tabcolsep}{3pt}
\renewcommand{\arraystretch}{1.05}

\begin{tabular}{lrrr}
\toprule
Dataset & Content sensitivity & Predictive utility & Reference--altered \\
\midrule
Bank & -.032 [-.044,-.019] & +.089 [+.078,+.101] & -.121 [-.136,-.106] \\
Blood & +.038 [-.020,+.095] & +.114 [+.032,+.197] & -.076 [-.156,+.003] \\
California & +.084 [+.070,+.097] & +.053 [+.039,+.066] & +.031 [+.018,+.042] \\
Car & +.409 [+.370,+.448] & +.300 [+.268,+.333] & +.108 [+.068,+.149] \\
Credit-g & +.066 [+.005,+.129] & -.003 [-.047,+.039] & +.069 [+.015,+.124] \\
Diabetes & +.087 [+.053,+.122] & +.164 [+.102,+.224] & -.077 [-.138,-.016] \\
Heart & -.010 [-.067,+.048] & -.021 [-.051,+.009] & +.011 [-.040,+.063] \\
Income & +.139 [+.132,+.146] & +.074 [+.068,+.079] & +.066 [+.057,+.074] \\
Jungle & +.230 [+.221,+.240] & +.025 [+.017,+.032] & +.206 [+.197,+.214] \\
\midrule
Macro & +.112 [+.036,+.206] & +.088 [+.032,+.157] & +.024 [-.040,+.091] \\
\bottomrule
\end{tabular}
\end{table}

At the macro level, both content sensitivity and predictive utility are positive.
At the dataset level, however, the two measures can differ substantially.
Credit-g has positive content sensitivity (\(+.066\)) but approximately zero predictive utility (\(-.003\)), while Bank shows negative content sensitivity (\(-.032\)) despite positive predictive utility (\(+.089\)).
These cases illustrate that the intended--altered difference does not determine the predictive benefit of the intended semantic information.

\subsection{Measurement knowledge}
\label{sec:app-measurement}

The focused three-endpoint analysis examines a setting in which coding polarity differs across schemas.
Without the intended decoding, reference performance falls below chance for diabetes and current smoking.
With the documented measurement information, both content sensitivity and predictive utility are positive.

Table~\ref{tab:app-measurement-full} also reports an extension to ten endpoints.
In this broader set, the reference AUROC is above chance on average, while predictive utility remains positive. The intended--altered comparison fails its frozen coverage requirement: cancer history in KNHANES-to-NHANES transfer has only 4 of 10 finite seed pairs, below the required eight. The frozen analysis retains available finite pairs within each endpoint--direction unit and weights endpoint clusters equally; this ten-endpoint comparison is descriptive and is not a coverage-qualified confirmation.

\begin{table}[h]
\caption{
Measurement-knowledge results with \(n_{\mathrm{target}}=256\).
Positive differences favor the intended condition.
The ten-endpoint content-sensitivity result is descriptive because its frozen finite-pair coverage gate failed.
}
\label{tab:app-measurement-full}
\centering
\footnotesize
\setlength{\tabcolsep}{3pt}

\begin{tabular}{
P{0.23\linewidth}
P{0.20\linewidth}
P{0.20\linewidth}
P{0.18\linewidth}
P{0.11\linewidth}}
\toprule
Panel
& Content sensitivity
& Predictive utility
& Reference--altered
& Reference AUROC \\
\midrule

Focused panel (3 endpoints)
& {+.292 [+.145, +.431]}
& {+.384 [+.202, +.553]}
& -.093 [-.228, +.012]
& .450 \\

Extended panel (10 endpoints)
& +.249 [+.182, +.309]
& {+.197 [+.110, +.294]}
& ---
& .556 [.499, .608] \\

\bottomrule
\end{tabular}
\end{table}

The focused panel shows positive content sensitivity and predictive utility under the polarity mismatch.
Predictive utility also remains positive in the ten-endpoint extension, although its magnitude is smaller.
The positive utility in the ten-endpoint extension shows that the result is not limited to the two below-chance reference endpoints, although the effect is smaller in the broader panel.

\subsection{Correspondence knowledge}
\label{sec:app-correspondence}

In the examination setting, the intended source--target correspondence outperforms a deranged correspondence, giving positive content sensitivity.
The original comparison, however, does not provide a suitable reference because removing the correspondence also changes the input representation.

We therefore construct a post-result width-matched reference with two columns allocated to each correspondence member in every condition.
In the intended and altered conditions, corresponding source and target values are placed in a shared member column, with the altered condition using a deranged source--target assignment.
In the reference condition, source and target values instead occupy distinct domain-local columns with complementary missingness.

This construction preserves the base features, source transfer, labeled-target and query rows, representation width, and one materialized member-value location per row before native missingness.
Feature subsampling is disabled in all three conditions.
The reference therefore removes shared cross-domain correspondence while preserving the remaining predictive information and representation width.

Under the same checks used in the audit, this construction satisfies P1--P6. No source or target value or row is removed, the same two columns are reserved for each correspondence member in every condition, and the model, preprocessing, and fitting procedure are unchanged.
The tested change is only whether corresponding source and target values share a cross-domain column.

\begin{table}[h]
\caption{
Examination correspondence results with \(n_{\mathrm{target}}=256\).
The width-matched reference preserves the input width and underlying values while removing the source--target correspondence.
}
\label{tab:app-exam-no-bridge}
\centering
\footnotesize
\setlength{\tabcolsep}{3.5pt}
\renewcommand{\arraystretch}{1.05}

\begin{tabular}{P{0.18\linewidth}P{0.46\linewidth}P{0.22\linewidth}}
\toprule
Quantity & Comparison & Mean [95\% CI] \\
\midrule

Content sensitivity
& Intended -- altered correspondence
& {+.070 [+.029, +.121]} \\

Predictive utility
& Intended -- width-matched reference
& {+.047 [+.017, +.084]} \\

Reference--altered
& Width-matched reference -- altered correspondence
& +.022 [+.002, +.043] \\

\bottomrule
\end{tabular}
\end{table}
Because this reference was constructed after the original correspondence result was known, its predictive-utility comparison is treated as a post-result reference analysis rather than as part of the original ablation design.
Both content sensitivity and predictive utility are positive in the six-endpoint examination panel, but the intended--altered difference is larger.
The difference between them reflects the lower performance of the altered correspondence relative to the reference.

We also evaluate the same comparison across 13 eligible examination endpoints.
Content sensitivity remains positive at \(+.057\) (95\% CI \([+.025,+.093]\)), whereas predictive utility is smaller and unresolved at \(+.014\) (95\% CI \([-.013,+.042]\)).
Thus, the broader endpoint set preserves the distinction between sensitivity to the supplied correspondence and its predictive benefit.


\paragraph{Performance under incorrect correspondences.}
To test whether an altered correspondence can itself reduce performance, we progressively replace intended bindings with plausible alternative matches.
AUROC loss increases as more incorrect bindings are introduced and approaches the loss under the full derangement
(Figure~\ref{fig:app-correspondence-decoy-damage}).

\begin{figure}[h]
\centering
\includegraphics[width=\linewidth]{figures/appendix_figure_a7_decoy_damage.pdf}
\caption{
AUROC loss as intended correspondences are replaced by plausible alternative bindings.
Thin lines show individual endpoints, the thick line shows their mean, and the dashed line shows the full correspondence derangement.
}
\label{fig:app-correspondence-decoy-damage}
\end{figure}

This result provides direct evidence that part of an intended--altered correspondence gap can arise from degradation under the altered condition.


\subsection{Relation knowledge}
\label{sec:app-relation}

Relation knowledge shows a large difference between content sensitivity and predictive utility.
Across the real-data transfer settings, the intended relations outperform reversed relations by a substantial margin, while predictive utility relative to the unconstrained reference remains close to zero.

To test whether this difference depends on how the altered condition is constructed, we also replace the intended relations with randomly selected relations while keeping the intended and reference conditions fixed.
The intended--random differences are much smaller than the intended--reversed differences.

\begin{table}[H]
\caption{
Relation-knowledge results under reversed and random altered conditions.
Content sensitivity compares the intended relations with the indicated altered condition, while predictive utility compares the intended relations with the unconstrained reference.
Positive values favor the intended condition.
}
\label{tab:app-relation-altered-construction}
\centering
\footnotesize
\setlength{\tabcolsep}{3pt}

\begin{tabular}{P{0.30\linewidth}P{0.22\linewidth}P{0.22\linewidth}P{0.20\linewidth}}
\toprule
Setting
& Reversed sensitivity
& Random sensitivity
& Predictive utility \\
\midrule

NHANES--KNHANES
& +.183 [+.081, +.298]
& -.001 [-.006, +.003]
& -.002 \\

NHANES/KNHANES--HRS
& +.266 [+.175, +.367]
& +.041 [+.016, +.077]
& +.002 \\

CAMELS-120
& +.144 [+.127, +.163]
& +.002 [+.000, +.003]
& .000 [-.001, +.001] \\

\bottomrule
\end{tabular}
\end{table}

The contrast is strongest for the reversed condition.
For NHANES--KNHANES and CAMELS-120, replacing reversal with random relations reduces content sensitivity to approximately zero, while predictive utility is unchanged because the intended and reference conditions are identical across the two comparisons.
For HRS, the random relations also reduce performance relative to the reference, producing a smaller but still positive intended--random difference.

The controlled benchmark shows the same pattern when the intended relations are known by construction.
For the additive setting with \(n_{\mathrm{target}}=32\), content sensitivity is \(+.304\) against reversed relations but \(+.126\) against random relations, while predictive utility is only \(+.024\).
Thus, the magnitude of content sensitivity depends strongly on the altered relation used for comparison and need not reflect the predictive utility of the intended relation.

\section{Sensitivity to the Choice of Reference}

We next examine whether predictive utility changes across multiple suitable references.
We use the TabLLM feature-name experiment because feature-name information can be removed while keeping the remaining input format unchanged.


\subsection{Evaluated TabLLM reference constructions}
\label{sec:app-tabllm-reference-audit}

Starting from the anonymous-name reference used in
Section~\ref{sec:app-tabllm-reproduction},
we construct a frozen library of 24 anonymous-identifier references.
For a dataset with $p$ fields, every reference uses the same set of $p$
zero-padded identifiers
(\texttt{feature\_001}, \texttt{feature\_002}, $\ldots$)
and differs only in the bijection between those identifiers and the fixed field lines.

The primary reference, denoted \texttt{ref\_00}, uses the original anonymous mapping.
The remaining 23 references use distinct non-identity permutations generated from a deterministic dataset-specific random stream.
Duplicate mappings are skipped, and the library size of 24 was fixed before scoring.

Across all candidates, the List Template, field order, values, placeholders, punctuation, and number of fields remain unchanged.
Each candidate also preserves a stable field-to-identifier mapping across rows.
Thus, the evaluated references differ only in which anonymous identifier is assigned to each feature.

For Blood, which contains four fields, the 24 references exhaust all $4!$ possible assignments.
For the other datasets, they form a fixed subset of the larger admissible assignment space.
Accordingly, the reported ranges characterize sensitivity over the evaluated library rather than bounds over all possible anonymous references.

Across the 24 evaluated references, mean predictive utility over the nine datasets ranges from
\(+.088\) to \(+.124\).
At the dataset level, the range is substantially wider and crosses zero for four of the nine datasets.

\begin{table}[h]
\caption{
Predictive utility across the 24 evaluated anonymous-name references.
Canonical denotes the primary anonymous reference used in Section~\ref{sec:app-tabllm-reproduction}.
Minimum and maximum are descriptive point estimates across the evaluated references.
}
\label{tab:app-tabllm-reference-envelope}
\centering
\footnotesize
\setlength{\tabcolsep}{3pt}

\begin{tabular}{@{}lrrrr@{}}
\toprule
Dataset & Canonical & Minimum & Maximum & Span \\
\midrule
Bank       & +.089 & +.022 & +.166 & .144 \\
Blood      & +.114 & +.093 & +.246 & .153 \\
California & +.053 & -.011 & +.105 & .116 \\
Car        & +.300 & +.270 & +.345 & .075 \\
Credit-g   & -.003 & -.040 & +.045 & .086 \\
Diabetes   & +.164 & +.088 & +.217 & .130 \\
Heart      & -.021 & -.021 & +.143 & .163 \\
Income     & +.074 & +.055 & +.106 & .051 \\
Jungle     & +.025 & -.020 & +.083 & .102 \\
\bottomrule
\end{tabular}
\end{table}

The point-estimate range crosses zero for California, Credit-g, Heart, and Jungle.
These ranges describe variation across the evaluated references; whether the apparent sign changes remain after accounting for evaluation uncertainty is examined in the next subsection.

\subsection{Reference variation and evaluation uncertainty}
\label{sec:app-tabllm-reference-noise}

The ranges in the previous subsection are based on point estimates.
We therefore test how much of the observed variation across references can be explained by evaluation uncertainty.

We use 5,000 paired bootstrap resamples over the shared test examples.
Within each resample, the same example weights are used for the intended condition and all 24 references, so uncertainty from the shared evaluation data is accounted for jointly.

Let $\hat{\mathbf u}$ denote the $R=24$ predictive-utility estimates for one dataset,
let $\widehat{\Sigma}$ denote their bootstrap covariance matrix, and define
\[
P = I-\mathbf{1}\mathbf{1}^{\top}/R.
\]
We measure the observed between-reference variance as
\[
V_{\mathrm{obs}}
=
\hat{\mathbf u}^{\top}P\hat{\mathbf u}/R,
\]
and estimate the contribution of evaluation noise to this spread as
\[
V_{\mathrm{noise}}
=
\operatorname{tr}(P\widehat{\Sigma})/R.
\]
The noise-corrected between-reference variance is therefore
$V_{\mathrm{obs}}-V_{\mathrm{noise}}$.
The projection removes common shifts shared across references, including uncertainty inherited from the common intended condition.

Across the nine datasets, 79.2--99.7\% of the observed between-reference variance remains after accounting for evaluation uncertainty.
Thus, most of the variation in utility across the evaluated references cannot be explained by test-set noise alone.

\begin{figure}[h]
\centering
\includegraphics[width=\linewidth]{figures/paired_reference_decomposition.pdf}
\caption{
Reference variation and evaluation uncertainty across the 24 evaluated TabLLM references.
\textbf{(a)} Observed between-reference variation, its estimate after accounting for evaluation uncertainty, and the typical uncertainty of an individual reference estimate.
\textbf{(b)} Simultaneous 95\% intervals for the references with the lowest and highest estimated utility in each dataset; gray points show the remaining reference estimates.
}
\label{fig:app-tabllm-reference-noise}
\end{figure}

The point-estimate range crosses zero for California, Credit-g, Heart, and Jungle.
After accounting for uncertainty across all 216 dataset--reference estimates, however, both a negative and a positive utility are resolved only for Jungle.
Its lowest estimated utility is \(-.0196\) (95\% CI \([-.0311,-.0081]\)), while its highest is \(+.0828\) (95\% CI \([+.0694,+.0962]\)).

For California, Credit-g, and Heart, at least one end of the range remains unresolved.
The 4/9 result therefore describes point-estimate sensitivity to reference choice, whereas a resolved change in sign is supported for 1/9 datasets.

\subsection{Additional reference checks}

\paragraph{Reproduction sensitivity.}
Because the released-code reproduction does not exactly match every published TabLLM result, we repeat the reference analysis on the four datasets for which the intended, permuted-name, and values-only conditions are all reproduced within .015 AUROC of the published means.
This restriction retains Car, Diabetes, Income, and Jungle and does not use the anonymous reference or predictive-utility results.

The main pattern remains in this subset.
Utility under the primary anonymous reference is \(+.141\), and the mean utility across the 24 evaluated references ranges from \(+.120\) to \(+.175\).
Dataset-level point-estimate ranges cross zero for one of the four datasets, compared with four of nine in the full panel.
Thus, positive aggregate utility and variation across references are not driven by the datasets with the largest reproduction differences.

\paragraph{Alternative anonymous labels.}
The primary reference uses artificial identifiers such as \texttt{feature\_001}.
To check whether the result depends on this particular wording, we repeat the zero-shot comparison using \texttt{Field A} and \texttt{Variable A} while keeping the feature mapping and the rest of the input format unchanged. Intervals for this fixed-mapping analysis use 2,000 paired row resamples, preserving overlap across the split seeds.

Mean predictive utility remains positive for all three versions: \(+.0887\) with \texttt{feature\_001}, \(+.0716\) with \texttt{Field A}, and \(+.0792\) with \texttt{Variable A}.
Relative to the primary reference, \texttt{Field A} reduces utility by \(-.0171\) (95\% CI \([-.0268,-.0077]\)), while the shift for \texttt{Variable A} is smaller at \(-.0095\).

These results show that the positive mean utility is not specific to the \texttt{feature\_001} labels, although its magnitude still depends on how the anonymous identifiers are written.


\section{Predictive Utility and Information Available to the Reference}

This section examines how predictive utility changes when the information supplied by semantic knowledge becomes easier or harder for the reference to recover.


\subsection{Measurement mismatch}
\label{sec:app-measurement-dose}

We vary the amount of measurement mismatch between the source and target schemas while keeping the target representation fixed.
Across all three outcome families and all evaluated models, predictive utility increases as more source features use their dataset-specific measurement representation
(Figure~\ref{fig:app-measurement-dose}).

With \(n_{\mathrm{target}}=32\), predictive utility at the largest mismatch is approximately \(+.26\) for XGB and HistGB and ranges from \(+.18\) to \(+.25\) for TabPFN.
With \(n_{\mathrm{target}}=512\), the corresponding effects are much smaller, ranging from \(+.02\) to \(+.04\).
For every model and outcome family, the change from \(n_{\mathrm{target}}=32\) to \(512\) is negative with a 95\% paired interval below zero.

\begin{figure}[h]
\centering
\includegraphics[width=\linewidth]{figures/appendix_figure_a2_mismatch_grid.pdf}
\caption{
Predictive utility as measurement mismatch increases.
Curves show the mean over 20 realizations for the additive, pairwise, and sparse outcome families.
The upper row uses \(n_{\mathrm{target}}=32\), and the lower row uses \(n_{\mathrm{target}}=512\).
}
\label{fig:app-measurement-dose}
\end{figure}

\begin{table}[h]
\caption{
Summary of the measurement-mismatch experiment across the three outcome families.
Ranges report the minimum and maximum family-level values.
}
\label{tab:app-measurement-dose-stats}
\centering
\footnotesize
\setlength{\tabcolsep}{4pt}

\begin{tabular}{lccc}
\toprule
Model & \(n_{\mathrm{target}}\) & Spearman \(\rho\) & Full-mismatch utility \\
\midrule
XGB    & 32  & .78--.90 & +.259--+.260 \\
HistGB & 32  & .79--.89 & +.257--+.264 \\
TabPFN & 32  & .75--.85 & +.176--+.249 \\
XGB    & 512 & .55--.77 & +.028--+.037 \\
HistGB & 512 & .63--.76 & +.030--+.035 \\
TabPFN & 512 & .57--.80 & +.022--+.035 \\
\bottomrule
\end{tabular}
\end{table}

The same qualitative pattern appears across outcome families and model classes: greater measurement mismatch increases the value of the supplied measurement information, while additional labeled target data substantially reduces that value.

\subsection{Relation reconstructibility}
\label{sec:app-relation-value-interface}
We distinguish two ways of supplying relation knowledge.
In the \emph{computed-value} interface, the model receives values obtained by applying the documented relation operation to its constituent variables.
In the \emph{executable-direction} interface, the relation instead constrains the admissible predictive function while the corresponding values remain otherwise available.

The reconstructibility analyses in this section concern the computed-value interface.
They test whether the utility of an explicitly supplied relation value depends on how easily that value can already be recovered from the remaining inputs.
They do not identify the utility of directional constraints:
for example, a relation value and its negation have the same reconstructibility from the underlying frame.
Directional-constraint utility is analyzed separately through the unconstrained-reference comparisons reported in
Section~\ref{sec:app-relation-capacity}.

We test whether predictive utility increases when the relation supplied to the model becomes harder to recover from the other input features.

When all constituent features are observed, adding the computed relation value provides almost no additional predictive utility.
We then progressively remove the constituent features used to compute that value while keeping the intended and reference conditions otherwise matched.
As the constituents are removed, predictive utility increases sharply.

For pairwise XGB with \(n_{\mathrm{target}}=32\), utility increases from \(+.002\) with all constituents present to \(+.242\) after all relevant constituents are removed.
The same pattern appears for the sparse family, both tree models, and both target-support settings.
Across these settings, the intended condition changes comparatively little, while reference performance deteriorates as the relation becomes harder to recover from the remaining inputs.

\begin{figure}[h]
\centering
\includegraphics[width=0.9\linewidth]{figures/appendix_figure_a8_redundancy_ladder.pdf}
\caption{
Performance as the constituent inputs of the supplied relation are progressively removed.
The intended condition changes comparatively little, whereas the reference deteriorates, increasing predictive utility.
Results are shown for the pairwise and sparse outcome families at two labeled-target sample sizes.
}
\label{fig:app-redundancy-ladder}
\end{figure}

This experiment changes both the available information and the number of input columns.
We therefore perform an additional fixed-width experiment in which the constituent features remain present but are progressively corrupted with noise.

\paragraph{Fixed-width reconstructibility check.}
\label{sec:app-noisy-proxy}

For the pairwise and sparse outcome families, each constituent feature remains present under its original identity but receives Gaussian noise at
\[
\sigma\in\{0,.125,.25,.5,1,2,4\},
\]
with noise scaled by the source interquartile range of that feature.
At every noise level, the intended and reference conditions receive the identical noisy eight-column base frame, including the same rows, missingness pattern, split, and learner.
Only the intended condition additionally receives the exact relation value computed from the clean constituent variables.

The experiment uses 20 realizations per outcome family, XGB and HistGradientBoosting, and
$n_{\mathrm{target}}\in\{32,512\}$.
The protocol was introduced after the earlier reconstructibility results were known but was fixed before this diagnostic was executed.
We therefore treat it as a post-result controlled diagnostic.
The reconstruction analysis was completed before predictive utility was evaluated.

As noise increases, the relation becomes harder to reconstruct from the remaining inputs and predictive utility increases.
For representative pairwise XGB with \(n_{\mathrm{target}}=32\), reconstruction \(R^2\) falls from .987 to .032 while utility increases from \(+.002\) to \(+.266\).
For the sparse family, reconstruction \(R^2\) falls from .988 to .146 while utility increases from approximately zero to \(+.231\).

The same pattern appears across all eight family--model--support settings.
At fixed nonzero noise levels, realizations with lower reconstruction $R^2$
generally have higher utility. At each dose, we first compute Spearman correlation between $1-R^2$ and utility across the 20 realizations separately in each family--learner--support setting, then take the arithmetic mean across the eight settings. These dose-specific means range from \(+.165\) to \(+.329\) across the six nonzero doses.
This within-dose analysis is post-result and descriptive.
Unlike the across-dose trend, it does not rely on the mechanically ordered noise dose and therefore provides complementary evidence for the same relationship without changing input width or variable identity.

\paragraph{Real-data boundary.}
\label{sec:app-relation-real-data}

We also examined whether the same relationship could be isolated in real cross-table transfer.
The available real-data cases did not provide a clean manipulation of reconstructibility: in CAMELS and NHANES--KNHANES, changing reconstructibility also changed other input information, and HRS contained no eligible computed-value relation.

We therefore treat the reconstructibility result as controlled-benchmark evidence rather than as a mechanism established in the real-data transfer settings.

\subsection{Labeled target support}

Across three settings, predictive utility decreases as more labeled target data become available.
We examine this pattern for TabLLM feature-name information, controlled measurement knowledge, and correspondence knowledge in examination transfer.

\paragraph{TabLLM.}
Using the primary anonymous reference, feature-name utility decreases from \(+.088\) at
\(n_{\mathrm{adapt}}=0\) to \(+.007\) at \(n_{\mathrm{adapt}}=512\).

\begin{table}[H]
\caption{
TabLLM feature-name predictive utility as labeled adaptation data increase.
Intervals are 95\% dataset-bootstrap intervals.
}
\label{tab:app-tabllm-support-ladder}
\centering
\footnotesize
\setlength{\tabcolsep}{4pt}

\begin{tabular}{lrr}
\toprule
\(n_{\mathrm{adapt}}\) & Predictive utility & 95\% CI \\
\midrule
0   & +.088 & [+.033,+.153] \\
4   & +.073 & [+.012,+.144] \\
32  & +.035 & [+.017,+.054] \\
512 & +.007 & [+.003,+.011] \\
\midrule
Decrease, \(0\rightarrow512\) & +.081 & [+.025,+.148] \\
\bottomrule
\end{tabular}
\end{table}

The decrease is not specific to the primary reference.
Across all 24 anonymous references evaluated in
Section~\ref{sec:app-tabllm-reference-audit},
utility ranges from \(+.088\) to \(+.124\) at zero shot and from \(+.002\) to \(+.010\) at
\(n_{\mathrm{adapt}}=512\).
The corresponding decreases range from \(+.081\) to \(+.116\), and the bootstrap interval for the smallest decrease across the 24 references is
\([+.017,+.138]\).

The same pattern remains with alternative anonymous labels.
At \(n_{\mathrm{adapt}}=512\), mean utility is \(+.0090\) with
\texttt{Field A} and \(+.0083\) with \texttt{Variable A}, compared with
\(+.0071\) for the primary \texttt{feature\_001} reference.

The aggregate decrease masks substantial variation across datasets.
Seven of nine datasets have lower utility at \(n_{\mathrm{adapt}}=512\) than at zero shot,
but only two trajectories decrease monotonically across all evaluated support levels.

\begin{figure}[H]
\centering
\includegraphics[width=0.82\linewidth]{figures/appendix_figure_tabllm_support_by_dataset.pdf}
\caption{
Dataset-level TabLLM feature-name utility across labeled adaptation levels.
Although mean utility decreases with additional labeled data, individual datasets show heterogeneous and often nonmonotonic trajectories.
}
\label{fig:app-tabllm-support-by-dataset}
\end{figure}


\paragraph{Controlled measurement knowledge.}
The controlled measurement experiment shows the same attenuation with labeled target support.
At the largest measurement mismatch, XGB predictive utility across the three outcome families decreases from \(+.259\)--\(+.260\) with
\(n_{\mathrm{target}}=32\) to \(+.028\)--\(+.037\) with
\(n_{\mathrm{target}}=512\).
The same qualitative pattern is observed for the other evaluated models
(Figure~\ref{fig:app-measurement-dose}).


\paragraph{Correspondence knowledge.}
In examination transfer, correspondence utility also decreases as the number of labeled target rows increases.
The six-endpoint subset was defined after the 13-endpoint expansion based on
whether the target was itself a matched variable, so we report the broader
13-endpoint analysis alongside it.

\begin{table}[h]
\caption{
Predictive utility of examination correspondence knowledge as labeled target data increase.
The six-endpoint panel contains targets that are themselves matched variables, while the 13-endpoint row reports the broader endpoint set.
}
\label{tab:app-exam-support-ladder}
\centering
\footnotesize
\setlength{\tabcolsep}{3pt}

\begin{tabular}{lrrrr}
\toprule
Endpoints
& \(n_{\mathrm{target}}=64\)
& \(n_{\mathrm{target}}=256\)
& \(n_{\mathrm{target}}=1024\)
& Decrease, \(64\rightarrow1024\) \\
\midrule

Six endpoints
& +.070 [+.037,+.102]
& +.047 [+.018,+.084]
& +.018 [+.008,+.028]
& +.052 [+.028,+.077] \\

All 13 endpoints
& +.039 [+.007,+.069]
& +.014 [-.013,+.042]
& +.008 [-.001,+.017]
& +.031 [+.003,+.056] \\

\bottomrule
\end{tabular}
\end{table}

In the six-endpoint panel, predictive utility decreases from \(+.070\) at
\(n_{\mathrm{target}}=64\) to \(+.018\) at \(n_{\mathrm{target}}=1024\).
The broader 13-endpoint analysis shows the same direction, although utility is unresolved at
\(n_{\mathrm{target}}=256\) and \(1024\).

Absolute performance improves in both conditions:
nRMSE decreases from .738 to .676 for the intended condition and from .808 to .694 for the reference.
The larger improvement of the reference accounts for the shrinking predictive utility.

The examination and controlled support manipulations differ in how labeled-target weight changes with sample size.
In the examination ladder, each labeled target row receives weight
$n_{\mathrm{source}}/n_{\mathrm{target}}$,
so the total target-domain training weight remains fixed while the number of distinct labeled target rows increases.
In the controlled benchmark, increasing $n_{\mathrm{target}}$ also increases the total target contribution because unit weights are used.
We therefore compare the two settings at the level of their qualitative attenuation pattern rather than treating the support levels as quantitatively equivalent interventions.

Across all three settings, additional labeled target information reduces the predictive utility of the supplied semantic knowledge.
The common pattern is that the reference improves as target information becomes more available, while the intended condition changes comparatively less.

\FloatBarrier

\section{Additional Robustness and Learner Coverage}


\subsection{Measurement robustness}

\paragraph{Polarity-invariant AUROC check.}
Because two reference AUROCs in the focused measurement panel are below chance,
we repeat the comparison using \(f(A)=\max(A,1-A)\).
Predictive utility remains positive in the three-endpoint panel,
\(+.276\) (95\% CI \([+.090,+.385]\)),
and in the ten-endpoint extension,
\(+.156\) (95\% CI \([+.089,+.225]\)).
Thus, the positive utility is not explained solely by AUROC polarity.

\paragraph{Rank-alignment check.}
The controlled measurement experiment combines differences in numerical representation with polarity reversals.
To separate these effects, we apply the same within-domain rank transformation to the intended and reference conditions before prediction.
This removes order-preserving differences in scale and coding while retaining polarity reversals.

With \(n_{\mathrm{target}}=32\), predictive utility remains positive in all six
outcome-family--model settings after rank alignment, with a mean of \(+.228\)
compared with \(+.260\) under the original representation.
With \(n_{\mathrm{target}}=512\), the remaining effect is much smaller,
ranging from \(+.008\) to \(+.019\) across the XGB outcome families.

After rank alignment, utility also increases with the number of reversed ordinal features.
At \(n_{\mathrm{target}}=32\), the mean increase is \(+.088\)
(95\% CI \([+.046,+.129]\)) per reversed feature.
At \(n_{\mathrm{target}}=512\), the corresponding slope decreases to \(+.004\)
(95\% CI \([+.001,+.008]\)).
When no ordinal feature is reversed, the rank-aligned intended and reference
inputs are identical, so predictive utility is exactly zero by construction.

\begin{table}[h]
\caption{
Summary of the controlled measurement robustness checks.
Rank-aligned utility removes order-preserving differences between the source
and target representations while retaining polarity reversals.
}
\label{tab:app-measurement-rank-alignment}
\centering
\footnotesize
\setlength{\tabcolsep}{4pt}

\begin{tabular}{lccc}
\toprule
Setting
& Original utility
& Rank-aligned utility
& Utility increase per reversal \\
\midrule

\(n_{\mathrm{target}}=32\), XGB
& +.259--+.260
& +.169--+.275
& +.088 [+.046,+.129] \\

\(n_{\mathrm{target}}=32\), HistGB
& +.257--+.264
& +.167--+.278
& +.088 [+.046,+.129] \\

\(n_{\mathrm{target}}=512\), XGB
& +.028--+.037
& +.008--+.019
& +.004 [+.001,+.008] \\

\bottomrule
\end{tabular}
\end{table}

Together, these checks show that the measurement result is not explained only by
below-chance AUROC or by order-preserving recoding.
After those effects are removed, the remaining utility is closely associated with
polarity mismatch and becomes much smaller when more labeled target data are available.


\subsection{Model and capacity dependence}

Relation utility varies with both model class and model capacity.
We first examine this dependence within the controlled benchmark and then extend the comparison to TransTab and CARTE.


\subsubsection{Tree capacity and MLP}
\label{sec:app-relation-capacity}

When the computed relation values are already available, adding the intended directional constraints provides a small positive increment across the tree-based settings.
Across the pairwise and sparse outcome families, this additional utility ranges from \(+.009\) to \(+.018\) over the two labeled-target sample sizes and the two tree models.

The magnitude of relation utility is sensitive to tree capacity.
Reducing XGBoost depth from 6 to 2 decreases utility by roughly an order of magnitude for the additive and pairwise families and makes the sparse-family mean slightly negative.

\begin{table}[H]
\caption{
Effect of XGBoost depth on relation predictive utility.
The baseline uses depth 6; the reduced-capacity model uses depth 2.
}
\label{tab:app-relation-capacity}
\centering
\footnotesize
\setlength{\tabcolsep}{4pt}

\begin{tabular}{lrrr}
\toprule
Family
& \(n_{\mathrm{target}}\)
& Depth 6
& Depth 2 \\
\midrule

Additive & 32  & +.0244 & +.0017 \\
Additive & 512 & +.0212 & +.0012 \\
Pairwise & 32  & +.0096 & +.0003 \\
Pairwise & 512 & +.0105 & +.0003 \\
Sparse   & 32  & +.0068 & -.0014 \\
Sparse   & 512 & +.0067 & -.0015 \\

\bottomrule
\end{tabular}
\end{table}

The MLP shows a different decomposition.
Predictive utility is positive for the pairwise and sparse families, ranging from \(+.006\) to \(+.013\), but remains unresolved for the additive family.
The altered directional constraints also behave differently across families: they reduce performance in the additive setting but can outperform the reference in the pairwise and sparse settings.
Thus, both predictive utility and the intended--altered difference depend on the model through which the relation information is used.


\subsubsection{TransTab and CARTE}
\label{sec:app-learner-coverage}

We further evaluate the same distinction with TransTab and CARTE.
The results are not uniformly positive across models, semantic components, or labeled-target sample sizes.

\paragraph{TransTab setup.}
We use the official TransTab implementation at commit
\texttt{fdb34cf38abda73ee6a741b802fe226cc89ba7b5}
with a \texttt{TransTabRegressor} initialized from scratch.
The model uses hidden dimension 128, two Transformer layers,
eight attention heads, a 256-dimensional feed-forward layer,
ReLU activations, and zero dropout.
We train for 50 fixed epochs with Adam, learning rate \(10^{-4}\),
zero weight decay, and batch size 64, without validation-based model selection.

Source and labeled-target tables are provided jointly through TransTab's
variable-column supervised interface.
Training labels are standardized using the mean and population standard
deviation of the source training labels and transformed back to the original
scale for evaluation.
Numeric missing values are filled using domain-local training medians with
explicit missingness indicators; target-query preprocessing reuses the state
fit on the labeled target-support data.
The realization index is used as the random seed and is held fixed across
paired semantic conditions.

Two compatibility corrections were required for regression execution:
we corrected the upstream \texttt{TrainDataset.\_\_getitem\_\_} slice from
\texttt{index-1:index} to \texttt{index:index+1}, and read the raw regression
logit rather than the classifier-sigmoid output returned by
\texttt{predict()}.
Neither change modifies the architecture, training objective, or optimizer.

\paragraph{CARTE setup.}
We use \texttt{carte-ai==0.0.26} with a single
\texttt{CARTEMultitableRegressor}.
The distributed pretrained weights and fastText embeddings are loaded,
with the pretrained representation weights kept frozen.
Models are fit with learning rate \(0.001\), batch size 64,
a maximum of 200 epochs, validation fraction \(0.2\),
early-stopping patience 20, and target fraction \(0.125\).
The realization index is used as the paired random state.

As for TransTab, numeric missing values are processed using domain-local
training medians and explicit missingness indicators, and the target-query
frame reuses preprocessing state fit on the target-support data.
Because CARTE does not provide the monotonic-constraint interface used by the
tree learners, relation knowledge is supplied through derived
value-plus-text columns rather than as an executable directional constraint.
The CARTE relation result should therefore be interpreted as learner coverage
for this CARTE-specific relation interface, not as a direct replication of
the tree constraint implementation.

The exact upstream CARTE Git revision was not recorded in the frozen artifact;
we therefore anchor reproducibility to
\texttt{carte-ai==0.0.26} and the hashes of the distributed pretrained and
fastText assets retained in the artifact.

\begin{table}[h]
\caption{
Additional model coverage.
For CARTE, counts report the number of outcome families with a positive predictive-utility interval out of three evaluated families.
}
\label{tab:app-model-coverage-compact}
\centering
\footnotesize
\setlength{\tabcolsep}{4pt}

\begin{tabular}{P{0.20\linewidth}P{0.24\linewidth}P{0.48\linewidth}}
\toprule
Model / support
& Component
& Result \\
\midrule

TransTab / \(n_{\mathrm{target}}=32\)
& Shared cross-domain identity
& Utility \(+.114\)--\(+.129\); 74--80\% of the intended--altered gap is associated with altered-condition degradation \\

TransTab / \(n_{\mathrm{target}}=512\)
& Shared cross-domain identity
& Utility \(+.007\)--\(+.009\); 97--100\% of the intended--altered gap is associated with altered-condition degradation \\

CARTE / \(n_{\mathrm{target}}=32\)
& Measurement / correspondence / relation
& Positive utility in 0/3, 0/3, and 0/3 outcome families, respectively \\

CARTE / \(n_{\mathrm{target}}=512\)
& Measurement / correspondence / relation
& Positive utility in 2/3, 2/3, and 0/3 outcome families, respectively \\

\bottomrule
\end{tabular}
\end{table}

For TransTab, shared anonymous identifiers outperform domain-specific anonymous identifiers, whereas intended lexical names do not show a resolved improvement over the shared identifiers.
Because the regression head is newly initialized, this comparison supports the contribution of shared cross-domain identity rather than pretrained lexical semantics.

For CARTE, positive utility appears for measurement and correspondence knowledge in some \(n_{\mathrm{target}}=512\) outcome families, while relation utility remains unresolved.
Together with the tree-capacity and MLP results, these analyses show that the measured predictive utility of semantic knowledge depends on the model and how that knowledge is incorporated.


\subsection{Fine-grained and unit-level analyses}
\label{sec:app-controlled-boundaries}

\paragraph{Correspondence interface check.}
Correspondence sensitivity depends on how the correspondence is supplied.
The intended--altered difference is \(+.065\) with explicit examination binding but approximately zero under canonical aggregation.
Survey binding gives positive differences with both XGB (\(+.045\)) and HistGB (\(+.051\)).
The survey no-added-columns comparison is reported only as an ablation because it removes additional input information and does not satisfy Preservation.

\paragraph{Individual CAMELS relations.}
The CAMELS content-sensitivity result is driven primarily by the precipitation direction.
Flipping precipitation increases log-RMSE by \(+.191\) (95\% CI \([+.143,+.248]\)), whereas flipping PET has approximately no effect,
\(-.001\) (95\% CI \([-.003,+.001]\)).
When both relations are considered together, the corresponding contributions are \(+.188\) for precipitation and \(-.003\) for PET.
Thus, the aggregate CAMELS result should not be interpreted as evidence that both declared relations contribute equally.

\paragraph{Resolution of unresolved effects.}
Zero-containing intervals differ substantially in precision.
For the clinical relation comparisons, the 80\% minimum detectable effect is
.0095--.0122 nRMSE, compared with .0015--.0029 log-RMSE for CAMELS.
The CAMELS-120 utility estimate is approximately zero with a narrow 95\% interval of
\([-.0011,+.0012]\) log-RMSE.
The three-endpoint survey correspondence comparisons have substantially lower resolution, with minimum detectable effects of .1443--.1586 AUROC.
Because these experiments use different metrics, the magnitudes are not compared directly, and the narrow CAMELS interval is not interpreted as formal equivalence to zero.

\paragraph{Unit-level heterogeneity.}
Aggregate results also hide variation across individual endpoints and basins.
Relation content sensitivity is positive across all evaluated HRS endpoints and the original CAMELS basins, whereas predictive utility varies around zero.
Survey correspondence results are more heterogeneous across endpoints.

\begin{table}[h]
\caption{
Unit-level variation in the main real-data comparisons.
Ranges summarize descriptive endpoint- or basin-level effects.
}
\label{tab:app-unit-attribution-summary}
\centering
\footnotesize
\setlength{\tabcolsep}{3pt}

\begin{tabular}{
P{0.27\linewidth}
P{0.14\linewidth}
P{0.26\linewidth}
P{0.26\linewidth}}
\toprule
Setting & Units & Content sensitivity & Utility / ablation \\
\midrule

Survey correspondence / XGB
& 3 endpoints
& all positive; +.022--+.072
& -.010--+.076 \\

Survey correspondence / HistGB
& 3 endpoints
& all positive; +.025--+.089
& -.024--+.073 \\

NHANES--KNHANES relation
& 6 endpoints
& all positive; +.010--+.445
& -.013--+.007 \\

NHANES/KNHANES--HRS relation
& 7 endpoints
& all positive; +.068--+.531
& -.014--+.019 \\

CAMELS relation
& 12 basins
& all positive; +.032--+.463
& -.004--+.008 \\

\bottomrule
\end{tabular}
\end{table}

\subsection{Exploratory cross-component interaction}
\label{sec:app-composition}

As an exploratory analysis, we examine whether measurement and correspondence
knowledge interact rather than contributing independently.
For measurement condition \(m\) and correspondence condition \(c\), we define

\[
I_{MC}
=
Q_{11}-Q_{10}-Q_{01}+Q_{00}.
\]

In the controlled benchmark, the interaction is positive at low labeled-target
support and decreases as source-label information is corrupted.
The initial real-data perturbations do not show the same attenuation pattern
(Figure~\ref{fig:app-interaction-controlled-real}).

\begin{figure}[h]
\centering
\includegraphics[width=\linewidth]
{figures/appendix_figure_controlled_real_mechanism.png}
\caption{
Exploratory measurement--correspondence interaction analysis.
\textbf{(a)} In the controlled benchmark, interaction decreases as source-label
information is corrupted.
\textbf{(b)} The initial real-data probes do not show the same attenuation pattern.
}
\label{fig:app-interaction-controlled-real}
\end{figure}

A later real-data check shows attenuation under labeled-target-label corruption
but not under source rendering or source-column permutation.
Because the result depends on the perturbation, we treat this analysis as
exploratory rather than evidence for a general cross-component mechanism.

\section{Published Semantic-Ablation Audit}
\label{sec:published_audit}

\subsection{Selection and coding procedure}

\paragraph{Study selection.}
We conduct a bounded retrospective audit of published semantic-ablation experiments rather than a systematic literature review.
The candidate pool consists of cross-table and semantic tabular-learning studies examined in our related-work review, together with closely related frameworks considered during revision.
A study is included if it reports at least one predictive comparison in which an identifiable semantic component supplied to the model is explicitly altered or removed.
Comparisons that vary pretraining, architecture, optimization, or representation mechanisms while leaving the evaluated semantic information unchanged do not satisfy this criterion.
A comparison can still enter the audit if it removes or alters the tested semantic component while also changing another part of the method; those additional changes are then evaluated under the criteria below.

For each included study, we retain all reported comparisons that satisfy this rule rather than selecting comparisons based on their effect direction or magnitude.
We additionally screened XTab and CM2 during revision.
Among the excluded studies, TransTab uses semantic column and categorical descriptions but does not report an ablation that alters or removes them; XTab does not use feature-level semantic information as a model input; and CM2 uses column-name semantics but likewise reports no ablation that alters or removes them.
Table~\ref{tab:app-audit-screening} summarizes the screening decisions. The artifact retains the original ten-study screening sheet and a separately dated finalization record of the author-verified revision decisions for XTab and CM2; the latter is not presented as a pre-revision record.
We screened 12 candidate studies and retained nine, yielding 25 qualifying comparisons.

\begin{table}[h]
\caption{Study screening for the published semantic-ablation audit.}
\label{tab:app-audit-screening}
\centering
\footnotesize
\setlength{\tabcolsep}{3pt}

\begin{tabular}{@{}P{0.20\linewidth}P{0.12\linewidth}P{0.62\linewidth}@{}}
\toprule
Study & Included & Basis for decision \\
\midrule
LIFT       & Yes & Explicit feature-name semantic manipulations \\
TransTab   & No & Uses column and categorical descriptions, but no reported ablation alters or removes this semantic information \\
TabLLM     & Yes & Explicit name/value semantic manipulations \\
PLATO      & Yes & Explicit knowledge-graph ablations \\
CARTE      & Yes & Explicit semantic-representation ablations \\
FeatLLM    & Yes & Explicit feature-description ablation \\
TabuLa-8B  & Yes & Explicit header-semantic comparison \\
ConTextTab & Yes & Explicit semantic-representation/name ablations \\
TabSTAR    & Yes & Explicit verbalization ablations \\
TARTE      & Yes & Explicit semantic-encoding comparison \\
XTab       & No & Does not use feature-level semantic information as a model input \\
CM2        & No & Uses column-name semantics, but no reported ablation alters or removes them \\
\bottomrule
\end{tabular}
\end{table}
Our post-result TransTab experiments are separate learner-coverage analyses and are not part of the published evidence used to determine this study's audit eligibility.

\paragraph{Reported claims.}
For each comparison, two human coders (one author and one independent researcher) independently examined the reported result, its caption, and the surrounding interpretation.
They recorded whether the study explicitly used the comparison to support a content-sensitivity claim, a predictive-utility claim, both, or neither.
Ambiguous cases were coded as unclear.

A content-sensitivity claim states that prediction depends on the particular semantic content supplied.
A general statement that semantic information is used or matters is not sufficient.
A predictive-utility claim states that the tested semantic knowledge improves predictive performance.

\paragraph{Control support.}
The coders then independently assessed what each reported control can establish.

For content sensitivity, the control must change the tested semantic content while preserving how that content is used and keeping the rest of the comparison unchanged.
For predictive utility, we apply the three criteria defined in the main paper: Removal, Preservation, and Interface.
Because Preservation had lower agreement in the initial coding, the same two coders then applied six fixed checks to all 25 comparisons before the final support labels were set. They asked whether (P1) the other predictive values were preserved, (P2) the feature set and feature identity were preserved, (P3) the input width was preserved, (P4) the rest of the input format was preserved, (P5) the model was unchanged, and (P6) the learning and inference procedure was unchanged. Each check was coded Yes, No, or Unclear from the paper and supplement, with the source passage recorded. Preservation is Yes only when P1--P4 are all Yes, and Interface is Yes only when P5--P6 are both Yes. If any required check is No, the criterion fails; otherwise, if a required check is Unclear, the criterion is Unclear. The six checks were fixed before the two coders applied them.

A predictive-utility control is coded as \emph{Supported} when all three criteria are satisfied, \emph{Unsupported} when at least one criterion clearly fails, and \emph{Indeterminate} when the published description does not provide enough information to establish all three criteria.

We then compare the reported claim with what its control supports.
An explicit claim is supported when the corresponding control meets its requirements, unsupported when a required condition clearly fails, and indeterminate when the available description is insufficient.
Comparisons that do not make the corresponding claim are not counted as unsupported claims, and unclear claim statements are kept separate.


\subsection{Coding agreement}

Two human coders independently coded all 25 comparisons before any discussion. Table~\ref{tab:app-published-audit-agreement} reports agreement for this initial coding.

\begin{table}[H]
\caption{
Agreement between the two human coders before adjudication.
}
\label{tab:app-published-audit-agreement}
\centering
\small
\setlength{\tabcolsep}{5pt}

\begin{tabular}{lcc}
\toprule
Coding item & Agreement & Cohen's \(\kappa\) \\
\midrule
Content-sensitivity claim   & 25/25 (100\%) & 1.000 \\
Predictive-utility claim    & 24/25 (96\%)  & .907 \\
Content-sensitivity support & 24/25 (96\%)  & .815 \\
Removal                     & 25/25 (100\%) & 1.000 \\
Preservation                & 15/25 (60\%)  & .265 \\
Interface                   & 22/25 (88\%)  & .784 \\
Predictive-utility support  & 19/25 (76\%)  & .649 \\
\bottomrule
\end{tabular}
\end{table}

Preservation had the lowest agreement in the initial coding, with agreement on 15 of 25 comparisons ($\kappa=.265$).
The same two coders then applied the fixed P1--P6 checks independently to all 25 comparisons and recorded the source passage for each answer. They agreed on 105 of the 150 individual checks (70.0%). Across the 18 explicit predictive-utility claims, the support labels derived from these independent checks agreed in 15 cases (83.3%, $\kappa=.693$).
After both coding sheets were complete, cases with different answers or questions raised in a consistency check were reviewed with the source passages, and final codes were set after agreement.
The final consensus codes are not treated as an inter-rater agreement measure. The full P1--P6 coding and source passages are included in the artifact.

AI tools were used only to organize cases and locate the relevant source passages. The coders made the judgments, and the authors reviewed the final coding.

\subsection{Full audit results}

Table~\ref{tab:app-published-audit} reports the comparison-level coding for all 25 audited ablations.
For each claim type, we show whether the original study made the claim and whether the reported control supports that claim under our criteria.

\begin{table}[H]
\caption{
Comparison-level coding for the 25 published semantic-ablation comparisons.
Claim is coded as Y (explicit claim), N (no corresponding claim), or ? (unclear).
Control support is coded as S (Supported), U (Unsupported), or I (Indeterminate);
a dash indicates that the corresponding claim was not made.
}
\label{tab:app-published-audit}
\centering
\scriptsize
\setlength{\tabcolsep}{3.5pt}
\renewcommand{\arraystretch}{1.02}

\begin{tabular}{@{}llcccc@{}}
\toprule
& &
\multicolumn{2}{c}{Content sensitivity} &
\multicolumn{2}{c}{Predictive utility} \\
\cmidrule(lr){3-4}\cmidrule(lr){5-6}
Study & ID & Claim & Support & Claim & Support \\
\midrule

LIFT       & C01-A & Y & S & Y & U \\
LIFT       & C01-B & Y & S & Y & U \\
LIFT       & C01-C & N & -- & Y & U \\
LIFT       & C01-D & N & -- & Y & I \\

\addlinespace
TabLLM     & C03-A & Y & S & N & -- \\
TabLLM     & C03-B & N & -- & N & -- \\
TabLLM     & C03-C & Y & U & N & -- \\

\addlinespace
PLATO      & C04-A & N & -- & Y & I \\
PLATO      & C04-B & N & -- & Y & U \\
PLATO      & C04-C & N & -- & N & -- \\

\addlinespace
CARTE      & C05-A & N & -- & Y & U \\
CARTE      & C05-B & N & -- & Y & I \\
CARTE      & C05-C & N & -- & Y & U \\

\addlinespace
FeatLLM    & C06-A & N & -- & Y & I \\
TabuLa-8B  & C07-A & ? & -- & Y & S \\

\addlinespace
ConTextTab & C08-A & N & -- & Y & U \\
ConTextTab & C08-B & N & -- & Y & U \\
ConTextTab & C08-C & N & -- & Y & U \\
ConTextTab & C08-D & N & -- & Y & U \\
ConTextTab & C08-E & N & -- & Y & I \\
ConTextTab & C08-F & N & -- & ? & I \\

\addlinespace
TabSTAR    & C09-A & N & -- & ? & S \\
TabSTAR    & C09-B & N & -- & Y & I \\

\addlinespace
TARTE      & C10-A & N & -- & N & -- \\
TARTE      & C10-B & N & -- & Y & U \\

\bottomrule
\end{tabular}
\end{table}

Four comparisons make an explicit content-sensitivity claim.
Three of these use controls that support the claim, while one uses a control that does not.

Eighteen comparisons make an explicit predictive-utility claim.
Of these, 1 uses a control that satisfies all three requirements, 11 use controls that fail at least one requirement, and 6 remain indeterminate because the published description does not establish all required properties.
Five additional comparisons make no predictive-utility claim, and two have unclear claim scope.

These counts describe the 25 comparisons included in this audit and are not intended to estimate the prevalence of unsupported controls in the broader literature.
An unsupported coding means that the reported control does not identify the stated claim under our criteria; it does not imply that the substantive claim is false or that the underlying method is invalid.

Preservation accounts for most of the coding disagreement reported above.
When the published description did not make clear whether other predictive information was preserved, we retained the Indeterminate category rather than assigning Supported or Unsupported.


\subsection{Detailed coding decisions}

Several cases required additional judgment about the scope of the reported claim
or about what information the control preserved.

For LIFT C01-C, we interpret the reported claim as a claim about the predictive
benefit of feature names.
The no-name prompt also removes task descriptions and categorical-value
verbalizations, so the comparison changes more than feature-name information.
We therefore code Preservation as not satisfied for this claim.

For TabLLM C03-B, \emph{List Only Values} removes feature names together with
the explicit key--value structure of the List Template.
Because our feature-name utility comparison requires the remaining input
structure and feature identity to be preserved, this control does not satisfy
Preservation.

For TabLLM C03-C, categorical values are relabeled, but continuous values are
first discretized into ten bins before permutation.
The comparison therefore changes numerical predictive information in addition
to the supplied value semantics.
We consequently do not treat it as isolating content sensitivity to value
semantics alone.

For CARTE C05-A, replacing the string feature initialization with MinHash changes the feature representation, so Preservation is not satisfied.

For FeatLLM C06-A, the evaluated component is the feature-description block supplied for rule generation.
Whether \emph{-Description} also removed the same descriptions from the rule-parsing prompt is not reported, so P4 is Unclear and the control is Indeterminate.

For ConTextTab C08-A--D, the ablations replace the semantic representations of
categorical or string features with ordinal, MinHash, Gap, or AutoGluon
encoders.
These comparisons therefore change the feature representation itself in
addition to removing the evaluated semantic information, so they do not
satisfy Preservation for a predictive-utility claim.
C08-E replaces column names with generic identifiers, but the published description does not establish whether the same learning procedure was used, so Interface is Unclear and the control is Indeterminate. 
The same issue applies to C08-F.

For TabSTAR C09-B, the paper does not state whether removing numerical information from the verbalization also changes the separate numerical input.
Preservation and Interface therefore remain Unclear, and the control is Indeterminate.

\subsection{Source evidence for coded claims}

Claim coding follows the interpretation stated by the original study rather
than the direction or magnitude of the reported ablation result.
For every comparison coded as making an explicit or unclear predictive-utility
claim, Table~\ref{tab:app-published-claim-provenance} records the interpretation
used for coding and the corresponding location in the original paper.
Comparisons for which we found no predictive-utility claim remain coded as
having no such claim.

\begin{table}[p]
\caption{
Source evidence used to code predictive-utility claims in the published-ablation audit.
The table includes comparisons with explicit or unclear predictive-utility claims.
Claim basis summarizes the interpretation given in the original study.
}
\label{tab:app-published-claim-provenance}
\centering
\scriptsize
\setlength{\tabcolsep}{3pt}
\renewcommand{\arraystretch}{1.04}

\begin{tabular}{@{}llcP{0.50\linewidth}P{0.20\linewidth}@{}}
\toprule
Study & ID & Claim & Claim basis & Source \\
\midrule

LIFT & C01-A & Y &
The gain over shuffled names is attributed to correct feature--value association in prompt format I. &
Sec.~4.1 / Table~9, p.~8 \\

LIFT & C01-B & Y &
The same interpretation attributes the gain in prompt format II to correct feature--value association. &
Sec.~4.1 / Table~9, p.~8 \\

LIFT & C01-C & Y &
Correctly incorporating feature names is stated to improve performance except on CMC; this interpretation also covers the no-name comparison. &
Sec.~4.1 / Table~9, p.~8 \\

LIFT & C01-D & Y &
The same qualified feature-name benefit is stated for the second no-name comparison. &
Sec.~4.1 / Table~9, p.~8 \\

PLATO & C04-A & Y &
Broader domain knowledge is interpreted as improving prediction beyond a feature-only knowledge graph. &
Sec.~4.1 / Table~3, pp.~7, 9 \\

PLATO & C04-B & Y &
Feature information in the knowledge graph is interpreted as improving prediction relative to the no-KG configuration. &
Sec.~4.1 / Table~3, pp.~7, 9 \\

CARTE & C05-A & Y &
Semantic-similarity encoding is described as important for exploiting external information, especially when the table itself provides little information. &
App.~C.3 / Fig.~10, p.~21 \\

CARTE & C05-B & Y &
Edge information is described as important for capturing table context and achieving strong predictive performance. &
App.~C.3 / Fig.~10, p.~21 \\

CARTE & C05-C & Y &
The attention layer is described as important for capturing table context and achieving strong predictive performance. &
App.~C.3 / Fig.~10, p.~21 \\

FeatLLM & C06-A & Y &
Feature descriptions and reasoning instructions are interpreted as using LLM prior knowledge to improve performance, especially in few-shot settings. &
Sec.~4.2 / Table~4, p.~8 \\

TabuLa-8B & C07-A & Y &
Meaningful headers are explicitly described as benefiting prediction at small shot counts, with the benefit decreasing as more labeled examples are available. &
App.~F.2 / Fig.~12, pp.~26--27 \\

ConTextTab & C08-A & Y &
The ordinal-encoding ablation is interpreted as losing performance because feature semantics are discarded. &
Sec.~5.1 / Table~2, pp.~8--9 \\

ConTextTab & C08-B & Y &
The MinHash ablation is interpreted as reducing the use of feature semantics. &
Sec.~5.1 / Table~2, pp.~8--9 \\

ConTextTab & C08-C & Y &
The same feature-semantics interpretation is applied to the Gap-encoding comparison. &
Sec.~5.1 / Table~2, pp.~8--9 \\

ConTextTab & C08-D & Y &
The same interpretation is applied to the AutoGluon-encoder comparison. &
Sec.~5.1 / Table~2, pp.~8--9 \\

ConTextTab & C08-E & Y &
The loss after replacing column names with generic identifiers is interpreted as evidence that the model benefits from column-name semantics. &
Sec.~5.1 / Table~2, pp.~8--9 \\

ConTextTab & C08-F & ? &
Enriched column descriptions are reported to provide a slight improvement, explicitly qualified as not statistically significant. &
Sec.~5.1 / Table~2, pp.~8--9 \\

TabSTAR & C09-A & ? &
Adding quantile information beyond bins is described as having limited impact despite a small average improvement. &
App.~G.4, p.~52 \\

TabSTAR & C09-B & Y &
Including numerical information in verbalization is interpreted as improving performance relative to feature names alone. &
Sec.~6 (Q3) / Table~4, pp.~9--10; App.~G.4, p.~52 \\

TARTE & C10-B & Y &
The MinHash comparison is interpreted as showing the importance of FastText semantic similarity for prediction. &
Sec.~4.3 / Fig.~6, pp.~9--10 (author version) \\

\bottomrule
\end{tabular}
\end{table}

\section{Limitations and Future Work}
\label{sec:app-limitations}

Our framework requires a reference that removes the semantic component being evaluated while preserving the rest of the predictive comparison.
Such a reference is not always available.
In some settings, removing semantic information can also change other predictive information or the model input, while in others the available description may not provide enough detail to verify what is preserved.
In these cases, we treat predictive utility as not separately evaluable rather than infer it from a control that changes additional factors.

Our experiments focus on settings in which the evaluated semantic information can be identified from feature names, dataset documentation, codebooks, known schema structure, or explicit relations.
The framework is less straightforward when the relevant information is distributed across opaque or high-dimensional representations and cannot be isolated as a well-defined component.
Future work could study how to construct and validate suitable references in such settings.

Predictive utility can also vary with the chosen reference and with how easily the learner can recover the removed information from other inputs or labeled target examples.
Our experiments examine this dependence across several learners and reference constructions, but they do not cover all possible models or representations of semantic knowledge.
Broader evaluations could test whether the same patterns hold for other tabular foundation models and cross-schema learning systems.

Finally, our published-ablation analysis is a bounded audit rather than a systematic review of the tabular-learning literature.
Its purpose is to examine whether specific claim--control pairs satisfy the proposed criteria, not to estimate how common particular evaluation practices are.
A larger systematic review could examine how often these patterns arise across a broader range of methods and domains.

\section{Reproducibility and Artifacts}
\label{sec:app-extraction}

\subsection{Numerical environment and access boundaries}

\begin{table}[h]
\caption{
Software versions used across the principal analyses and model-specific extensions.
The principal environment and recorded experiment-specific exceptions are listed separately; these version records do not imply that a complete historical environment lock exists for every run.
}\label{tab:app-software-environment}
\centering\small
\begin{tabular}{P{0.59\linewidth}P{0.33\linewidth}}
\toprule
Component & Version \\
\midrule
Python (principal analyses) & 3.10.19 \\
NumPy & 2.2.6 \\
pandas & 2.3.3 \\
scikit-learn & 1.7.2 \\
XGBoost & 3.1.3 \\
TabPFN & 7.0.0 \\
PyTorch & \texttt{2.12.0.dev20260408+cu128} \\
Transformers (documentation extraction) & 5.12.0 \\
Python (TabLLM / learner extensions) & 3.10.20 \\
Transformers (TabLLM / TransTab conformance) & 4.30.0 \\
Datasets (TabLLM) & 2.14.7 \\
TransTab & \texttt{fdb34cf38abda73ee6a74}\newline\texttt{1b802fe226cc89ba7b5} \\
CARTE & \texttt{carte-ai==0.0.26} \\
Python (examination support/reference, value-redundancy, CARTE) & 3.10.20 \\
Transformers (original TransTab run) & 5.12.0 \\
Relation-capacity sensitivity: Python / NumPy / XGBoost & 3.8.8 / 1.24.4 / 2.1.4 \\
Null-source extension: NumPy / scikit-learn / XGBoost & 1.24.4 / 1.3.2 / 2.1.4 \\

\bottomrule
\end{tabular}
\end{table}

The supplied one-thread replication reproduces all 400 comparable stored HistGradientBoosting nMSE values exactly. This is equality of recorded numerical outputs; cross-thread bitwise identity is not claimed. The artifact distinguishes exact equality from the $10^{-12}$ tolerance used to count material cross-thread differences, and documents the historical comparison scope.
The planned second HRS factorial pair was not executed because it requires new authorization for licensed HRS microdata.

The lightweight workflow verifies the paper from included aggregate outputs without restricted data or model calls. Full experiment execution additionally requires the documented source products, prepared inputs, model assets, and reviewer-specific access permissions. The artifact identifies incomplete historical environment locks and external clinical preparation dependencies explicitly; it does not claim that restricted-data experiments can be rerun from the archive alone.

\subsection{Supplementary artifact map}

\begin{table}[h]
\caption{Map of supplementary artifacts included with this submission. Public-data and governed-data execution dependencies are documented separately in the artifact.}
\label{tab:app-artifact-map}
\centering\small
\setlength{\tabcolsep}{4pt}
\begin{tabular}{P{0.30\linewidth}P{0.62\linewidth}}
\toprule
Artifact directory & Contents \\
\midrule
\path{provenance/} &
analysis-provenance ledger, with unavailable historical dates explicitly marked unknown \\
\path{tabllm/} &
complete frozen reference library, global-construction audit, product-space sensitivity summaries, and dataset-level reference envelopes
, exhaustive $k=512$ global-reference extension, and per-reference decay summaries \\
\path{robustness/measurement/} &
focused absolute-AUROC measurement display \\
\path{robustness/correspondence/} &
correspondence-interface display and clean-rerun record \\
\path{robustness/relation/} &
relation evidence audit, constituent-removal results, and full reconstructibility diagnostics \\
\path{robustness/interaction/} &
complete interaction surfaces, endpoint robustness, corruption calibration, and source-integrity analyses \\
\path{robustness/resolution/} &
full MDE calculations and unit-resolved results \\
\path{model_coverage/} &
complete TransTab and CARTE results \\
\path{method_contract/} &
complete prompt components, rendered prompts, schemas, acceptance rules, raw responses, and generation manifests \\
\path{published_audit/} &
initial human coding, independent P1--P6 checks, final consensus, source evidence, and adjudication history \\
\path{extraction/} &
complete seven-model measurement-code extraction panel \\
\path{reproducibility/} &
recorded software/model revisions, experiment-specific environment exceptions, missing-lock disclosures, and thread-count audit \\
\bottomrule
\end{tabular}
\end{table}


\FloatBarrier

\ICLRFinishAppendix
